% Complete reviewed Unit 8; editorial master: chapters/ch08/student.md
% Markdown SHA-256: 1b0ba8a7bc9fd6f7ad3b64c6958c45ecb7486fe8618eba82cc005cac48f14ba1
% One-use initial assembly followed by complete content and page review.
\newcommand{\runningtitle}{Unit 8. Questions and relatives}
% Layout inherited from the reviewed previous course; set a running title per unit.
% This is a small hand-maintained TeX layout, not a conversion pipeline.
\RequirePackage{fix-cm}
\documentclass[11pt,twoside]{article}
\usepackage{fontspec}
\setmainfont{Cambria}[
  Renderer=HarfBuzz,
  BoldFont=cambriab.ttf,
  ItalicFont=cambriai.ttf,
  BoldItalicFont=cambriaz.ttf
]
\setmonofont{lmmono10-regular.otf}[BoldFont=lmmonolt10-bold.otf]
\usepackage[paperwidth=165mm,paperheight=241mm,inner=21mm,outer=19mm,
  top=22mm,bottom=20mm,headheight=14pt,headsep=7mm,footskip=10mm]{geometry}
\usepackage{xcolor,array,booktabs,longtable,tabularx}
\usepackage{enumitem,fancyhdr}
\usepackage[unicode,hidelinks]{hyperref}
\definecolor{ink}{HTML}{000000}
\definecolor{muted}{HTML}{4C5558}
\setlength{\parindent}{0pt}
\setlength{\parskip}{3.5pt}
\AtBeginDocument{\fontsize{11}{14.3}\selectfont}
\setlength{\tabcolsep}{4pt}
\renewcommand{\arraystretch}{1.10}
\setlist[itemize]{leftmargin=1.3em,label=\textbullet,itemsep=2pt,topsep=3pt,parsep=0pt}
\setlist[enumerate]{leftmargin=1.8em,itemsep=3pt,topsep=4pt,parsep=0pt}
\widowpenalty=10000
\clubpenalty=10000
\emergencystretch=1em
\raggedbottom
\pagestyle{fancy}
\fancyhf{}
\providecommand{\runningtitle}{Old English}
\fancyhead[LE,RO]{\small\itshape \runningtitle}
\fancyhead[LO,RE]{\small Old English}
\fancyfoot[LE,RO]{\small\thepage}
\renewcommand{\headrulewidth}{0.3pt}
\newcommand{\booktitle}[2]{
  {\fontsize{9}{11}\selectfont\color{muted}#1\par}
  \vspace{3pt}
  {\fontsize{20}{24}\selectfont\color{ink}#2\par}
  \vspace{6pt}
}
\newcommand{\source}[1]{\par{\footnotesize\color{muted}#1\par}}
\newcommand{\readingtext}[1]{
  \par\addvspace{5pt}
  \begingroup\leftskip=1.5em
    \fontsize{12}{16}\selectfont #1\par
  \endgroup\addvspace{5pt}
}
\makeatletter
\renewcommand{\section}{\@startsection{section}{1}{\z@}
  {-2.4ex \@plus -0.6ex \@minus -0.2ex}{0.8ex \@plus 0.2ex}
  {\normalfont\fontsize{12}{15}\selectfont\bfseries\color{ink}}}
\renewcommand{\subsection}{\@startsection{subsection}{2}{\z@}
  {-1.8ex \@plus -0.5ex \@minus -0.2ex}{0.5ex \@plus 0.2ex}
  {\normalfont\fontsize{11}{14}\selectfont\bfseries\color{ink}}}
\makeatother

\hyphenpenalty=10000
\exhyphenpenalty=10000
\emergencystretch=3em
\newcommand{\itempara}[2]{\par\addvspace{3pt}\noindent\hangindent=1.8em\hangafter=1\makebox[1.8em][l]{#1.}#2\par}
\newcommand{\reserve}[1]{\par\penalty-100\begingroup\dimen0=\pagegoal\advance\dimen0 by-\pagetotal\ifdim\dimen0<#1\newpage\fi\endgroup}
\begin{document}
\addfontfeatures{Ligatures=TeXOff}
\raggedright
\booktitle{UNIT 8}{Questions and relative clauses}

% mdparagraph:P001:begin
This unit gives direct questions, the remaining selected pronoun and pronominal-adjective paradigms, and ordinary relative clauses. The additional verb systems and lexical constructions prepare the complete Cook I.8 exercise and two connected readings.
% mdparagraph:P001:end

% mdparagraph:P002:begin
Tables distinguish singular and plural, masculine, neuter and feminine, and nominative, genitive, dative and accusative. In verb tables, \textbf{1}, \textbf{2} and \textbf{3} mean first, second and third person. A row labelled “persons 1, 2, 3” includes every person in that number. Dual subjects take plural finite forms. Forms labelled \textbf{reference} are for recognition; their display does not require independent use of a new construction. Preserve vowel length and the distinction between e and ę.
% mdparagraph:P002:end

% mdparagraph:P003:begin
The examples and “The ship and the fish” are new teaching material. The merchant dialogue is an exact source excerpt with separate local help and original notes.
% mdparagraph:P003:end

\reserve{7\baselineskip}
% mdparagraph:P004:begin
\section*{1. Direct questions}
% mdparagraph:P004:end

% teaching: NEW-U08-GRAM-u08-direct-questions
% teaching: NEW-U08-GRAM-u08-clause-order
% mdparagraph:P005:begin
A direct \textbf{yes/\allowbreak{}no question} commonly begins with its finite verb, followed by the subject. The verb retains the person, number and tense required by that subject. Old English does not need an added verb corresponding to the English question auxiliary “do”.
% mdparagraph:P005:end

\reserve{8\baselineskip}
% example: NEW-U08-EXAMPLE-01
% mdparagraph:P006:begin
Hæfð se mann scip?
% mdparagraph:P006:end

% mdparagraph:P007:begin
Does the man have a ship?
% mdparagraph:P007:end

% mdparagraph:P008:begin
A question word or phrase commonly stands early, often before the finite verb. A question word that is itself the subject needs no second subject. These are common orders, not an exceptionless rule requiring one fixed arrangement. A phrase such as “whose ship” is one constituent, even though it contains two words.
% mdparagraph:P008:end

% mdparagraph:P009:begin
More generally, the finite verb commonly stands early in an independent clause and later, often finally, in a dependent clause. The relative clauses below illustrate that tendency. The case endings and the meaning of the clause remain essential when the order varies. This unit uses direct questions; indirect questions and general content complements come later.
% mdparagraph:P009:end

\reserve{7\baselineskip}
% mdparagraph:P010:begin
\subsection*{Who? What?}
% mdparagraph:P010:end

% teaching: NEW-U08-VOC-hwa
% teaching: NEW-U08-GRAM-u08-hwa-paradigm
% teaching: NEW-U08-GRAM-u08-interrogative-case
% mdparagraph:P011:begin
\textbf{Hwā} asks about a person; \textbf{hwæt} asks about a thing. Their paradigm is singular only. Masculine and feminine share one personal series; there is no separate feminine or plural paradigm. Choose the case from the interrogative’s own function: nominative subject, genitive possessor, accusative direct object, or dative recipient or other governed complement.
% mdparagraph:P011:end

% teaching: NEW-U08-PAR-HWA
\reserve{7\baselineskip}
% mdparagraph:P012:begin
\subsection*{hwā /\allowbreak{} hwæt}
% mdparagraph:P012:end

% table: NEW-U08-PAR-HWA
% mdtable:T01:begin
\par\begingroup
\fontsize{10.5}{13.2}\selectfont
\setlength{\LTleft}{0pt}\setlength{\LTright}{\fill}
\noindent
\begin{tabular}{@{}>{\raggedright\arraybackslash}p{25.54mm}>{\raggedright\arraybackslash}p{64.36mm}>{\raggedright\arraybackslash}p{28.48mm}@{}}
\toprule
\textbf{Case} & \textbf{Common masculine/\allowbreak{}feminine singular} & \textbf{Neuter singular} \\
\midrule
Nominative & hwā & hwæt \\
Genitive & hwæs & hwæs \\
Dative & hwǣm & hwǣm \\
Accusative & hwone & hwæt \\
\bottomrule
\end{tabular}
\par\endgroup

% mdtable:T01:end
\reserve{8\baselineskip}
% example: NEW-U08-EXAMPLE-02
% mdparagraph:P013:begin
Hwā hæfð ðæt scip? Hwone siehð se mann?
% mdparagraph:P013:end

% mdparagraph:P014:begin
Who has that ship? Whom does the man see?
% mdparagraph:P014:end

% mdparagraph:P015:begin
Hwā is the subject of hæfð. Hwone is the object of siehð; se mann remains the nominative subject. English “whom” corresponds here to an accusative because “see” has an accusative object.
% mdparagraph:P015:end

\reserve{8\baselineskip}
% example: NEW-U08-EXAMPLE-03
% mdparagraph:P016:begin
Hwǣm sęlst ðū ðone fisc? Hwæs scip siehst ðū?
% mdparagraph:P016:end

% mdparagraph:P017:begin
To whom do you give that fish? Whose ship do you see?
% mdparagraph:P017:end

% mdparagraph:P018:begin
Hwǣm is the dative recipient of the fish. Hwæs is genitive because it asks whose ship is seen; scip itself is the accusative object of siehst. The possessor and the possessed noun have different functions within that phrase.
% mdparagraph:P018:end

% mdparagraph:P019:begin
\textbf{Hwǣm} is the selected dative in this unit. Cook also gives \textbf{hwām} as a valid alternative. The distinct instrumental forms below are reference forms; learning them does not introduce general instrumental syntax.
% mdparagraph:P019:end

\reserve{7\baselineskip}
% mdparagraph:P020:begin
\subsection*{hwā reference variants and instrumentals}
% mdparagraph:P020:end

% table: NEW-U08-MATRIX-HWA-REFERENCE
% mdtable:T02:begin
\par\begingroup
\fontsize{10.5}{13.2}\selectfont
\setlength{\LTleft}{0pt}\setlength{\LTright}{\fill}
\noindent
\begin{tabular}{@{}>{\raggedright\arraybackslash}p{87.48mm}>{\raggedright\arraybackslash}p{33.71mm}@{}}
\toprule
\textbf{Function} & \textbf{Form} \\
\midrule
Dative common personal: printed alternative & hwām \\
Dative neuter: printed alternative & hwām \\
Neuter instrumental & hwȳ \\
Neuter instrumental alternative & hwon \\
\bottomrule
\end{tabular}
\par\endgroup

% mdtable:T02:end
\reserve{7\baselineskip}
% mdparagraph:P021:begin
\subsection*{How? Why?}
% mdparagraph:P021:end

% teaching: NEW-U08-VOC-hu
% teaching: NEW-U08-VOC-hwy
% teaching: NEW-U08-GRAM-u08-interrogative-adverbs
% mdparagraph:P022:begin
\textbf{Hū} asks about manner; \textbf{hwȳ} asks for a reason. Both are invariant question words in the following use. Lexical hwȳ “why?” can be used without constructing a general instrumental phrase.
% mdparagraph:P022:end

\reserve{7\baselineskip}
% mdparagraph:P023:begin
\subsection*{Question adverbs}
% mdparagraph:P023:end

% table: NEW-U08-MATRIX-QUESTION-ADVERBS
% mdtable:T03:begin
\par\begingroup
\fontsize{10.5}{13.2}\selectfont
\setlength{\LTleft}{0pt}\setlength{\LTright}{\fill}
\noindent
\begin{tabular}{@{}>{\raggedright\arraybackslash}p{60.59mm}>{\raggedright\arraybackslash}p{60.59mm}@{}}
\toprule
\textbf{Function} & \textbf{Form} \\
\midrule
How & hū \\
Why & hwȳ \\
\bottomrule
\end{tabular}
\par\endgroup

% mdtable:T03:end
\reserve{8\baselineskip}
% example: NEW-U08-EXAMPLE-04
% mdparagraph:P024:begin
Hwæt hæfð se mann? Hū swimð hē? Hwȳ swimð hē?
% mdparagraph:P024:end

% mdparagraph:P025:begin
What does the man have? How does he swim? Why does he swim?
% mdparagraph:P025:end

% mdparagraph:P026:begin
Hwæt is the accusative thing possessed. Hū and hwȳ ask different questions about the same swimming: how it is done, and why it is done.
% mdparagraph:P026:end

\reserve{7\baselineskip}
% mdparagraph:P027:begin
\section*{2. Interrogative adjectives}
% mdparagraph:P027:end

% teaching: NEW-U08-GRAM-u08-interrogative-adjectives
% mdparagraph:P028:begin
\textbf{Hwilc} means “which?” or “of what kind?”; \textbf{hwæðer}, “which of two?”; \textbf{hūlic}, “of what sort?”. They have strong adjective inflection and agree with their noun or referent in gender, number and case. No weak interrogative-adjective paradigm is needed here.
% mdparagraph:P028:end

% mdparagraph:P029:begin
The strong endings are the familiar adjective endings: masculine and neuter singular genitive -es and dative -um, masculine accusative -ne, feminine genitive and dative -re; plural genitive -ra and dative -um. The tables give the complete forms, including the nominative and accusative distinctions. Feminine plural -a is selected consistently; the previously noted feminine plural alternative -e remains valid.
% mdparagraph:P029:end

% teaching: NEW-U08-VOC-hwilc
% teaching: NEW-U08-PAR-HWILC-STRONG
\reserve{7\baselineskip}
% mdparagraph:P030:begin
\subsection*{hwilc: strong declension, singular}
% mdparagraph:P030:end

% table: NEW-U08-PAR-HWILC-STRONG-SG
% mdtable:T04:begin
\par\begingroup
\fontsize{10.5}{13.2}\selectfont
\setlength{\LTleft}{0pt}\setlength{\LTright}{\fill}
\noindent
\begin{tabular}{@{}>{\raggedright\arraybackslash}p{28.89mm}>{\raggedright\arraybackslash}p{28.89mm}>{\raggedright\arraybackslash}p{28.89mm}>{\raggedright\arraybackslash}p{28.89mm}@{}}
\toprule
\textbf{Case} & \textbf{Masculine} & \textbf{Neuter} & \textbf{Feminine} \\
\midrule
Nominative & hwilc & hwilc & hwilc \\
Genitive & hwilces & hwilces & hwilcre \\
Dative & hwilcum & hwilcum & hwilcre \\
Accusative & hwilcne & hwilc & hwilce \\
\bottomrule
\end{tabular}
\par\endgroup

% mdtable:T04:end
\reserve{7\baselineskip}
% mdparagraph:P031:begin
\subsection*{hwilc: strong declension, plural}
% mdparagraph:P031:end

% table: NEW-U08-PAR-HWILC-STRONG-PL
% mdtable:T05:begin
\par\begingroup
\fontsize{10.5}{13.2}\selectfont
\setlength{\LTleft}{0pt}\setlength{\LTright}{\fill}
\noindent
\begin{tabular}{@{}>{\raggedright\arraybackslash}p{28.89mm}>{\raggedright\arraybackslash}p{28.89mm}>{\raggedright\arraybackslash}p{28.89mm}>{\raggedright\arraybackslash}p{28.89mm}@{}}
\toprule
\textbf{Case} & \textbf{Masculine} & \textbf{Neuter} & \textbf{Feminine} \\
\midrule
Nominative & hwilce & hwilc & hwilca \\
Genitive & hwilcra & hwilcra & hwilcra \\
Dative & hwilcum & hwilcum & hwilcum \\
Accusative & hwilce & hwilc & hwilca \\
\bottomrule
\end{tabular}
\par\endgroup

% mdtable:T05:end
% mdparagraph:P032:begin
Keep the stem hwilc- before these endings. Hwelc and hwylc are documented spellings; expanded forms such as hwilcere and hwilcera also occur. The tables give the selected forms for recall.
% mdparagraph:P032:end

% teaching: NEW-U08-VOC-hulic
% teaching: NEW-U08-PAR-HULIC-STRONG
\reserve{7\baselineskip}
% mdparagraph:P033:begin
\subsection*{hūlic: strong declension, singular}
% mdparagraph:P033:end

% table: NEW-U08-PAR-HULIC-STRONG-SG
% mdtable:T06:begin
\par\begingroup
\fontsize{10.5}{13.2}\selectfont
\setlength{\LTleft}{0pt}\setlength{\LTright}{\fill}
\noindent
\begin{tabular}{@{}>{\raggedright\arraybackslash}p{28.89mm}>{\raggedright\arraybackslash}p{28.89mm}>{\raggedright\arraybackslash}p{28.89mm}>{\raggedright\arraybackslash}p{28.89mm}@{}}
\toprule
\textbf{Case} & \textbf{Masculine} & \textbf{Neuter} & \textbf{Feminine} \\
\midrule
Nominative & hūlic & hūlic & hūlic \\
Genitive & hūlices & hūlices & hūlicre \\
Dative & hūlicum & hūlicum & hūlicre \\
Accusative & hūlicne & hūlic & hūlice \\
\bottomrule
\end{tabular}
\par\endgroup

% mdtable:T06:end
\reserve{7\baselineskip}
% mdparagraph:P034:begin
\subsection*{hūlic: strong declension, plural}
% mdparagraph:P034:end

% table: NEW-U08-PAR-HULIC-STRONG-PL
% mdtable:T07:begin
\par\begingroup
\fontsize{10.5}{13.2}\selectfont
\setlength{\LTleft}{0pt}\setlength{\LTright}{\fill}
\noindent
\begin{tabular}{@{}>{\raggedright\arraybackslash}p{28.89mm}>{\raggedright\arraybackslash}p{28.89mm}>{\raggedright\arraybackslash}p{28.89mm}>{\raggedright\arraybackslash}p{28.89mm}@{}}
\toprule
\textbf{Case} & \textbf{Masculine} & \textbf{Neuter} & \textbf{Feminine} \\
\midrule
Nominative & hūlice & hūlic & hūlica \\
Genitive & hūlicra & hūlicra & hūlicra \\
Dative & hūlicum & hūlicum & hūlicum \\
Accusative & hūlice & hūlic & hūlica \\
\bottomrule
\end{tabular}
\par\endgroup

% mdtable:T07:end
% mdparagraph:P035:begin
Hūlic keeps hūlic- throughout. Cook marks the first vowel long; this chapter retains that notation. The separately documented spelling hūlīc does not change the selected table.
% mdparagraph:P035:end

% teaching: NEW-U08-VOC-hwaether
% teaching: NEW-U08-PAR-HWAETHER-STRONG
\reserve{7\baselineskip}
% mdparagraph:P036:begin
\subsection*{hwæðer: strong declension, singular}
% mdparagraph:P036:end

% table: NEW-U08-PAR-HWAETHER-STRONG-SG
% mdtable:T08:begin
\par\begingroup
\fontsize{10.5}{13.2}\selectfont
\setlength{\LTleft}{0pt}\setlength{\LTright}{\fill}
\noindent
\begin{tabular}{@{}>{\raggedright\arraybackslash}p{28.89mm}>{\raggedright\arraybackslash}p{28.89mm}>{\raggedright\arraybackslash}p{28.89mm}>{\raggedright\arraybackslash}p{28.89mm}@{}}
\toprule
\textbf{Case} & \textbf{Masculine} & \textbf{Neuter} & \textbf{Feminine} \\
\midrule
Nominative & hwæðer & hwæðer & hwæðer \\
Genitive & hwæðres & hwæðres & hwæðerre \\
Dative & hwæðrum & hwæðrum & hwæðerre \\
Accusative & hwæðerne & hwæðer & hwæðre \\
\bottomrule
\end{tabular}
\par\endgroup

% mdtable:T08:end
\reserve{7\baselineskip}
% mdparagraph:P037:begin
\subsection*{hwæðer: strong declension, plural}
% mdparagraph:P037:end

% table: NEW-U08-PAR-HWAETHER-STRONG-PL
% mdtable:T09:begin
\par\begingroup
\fontsize{10.5}{13.2}\selectfont
\setlength{\LTleft}{0pt}\setlength{\LTright}{\fill}
\noindent
\begin{tabular}{@{}>{\raggedright\arraybackslash}p{28.89mm}>{\raggedright\arraybackslash}p{28.89mm}>{\raggedright\arraybackslash}p{28.89mm}>{\raggedright\arraybackslash}p{28.89mm}@{}}
\toprule
\textbf{Case} & \textbf{Masculine} & \textbf{Neuter} & \textbf{Feminine} \\
\midrule
Nominative & hwæðre & hwæðer & hwæðra \\
Genitive & hwæðerra & hwæðerra & hwæðerra \\
Dative & hwæðrum & hwæðrum & hwæðrum \\
Accusative & hwæðre & hwæðer & hwæðra \\
\bottomrule
\end{tabular}
\par\endgroup

% mdtable:T09:end
% mdparagraph:P038:begin
Hwæðer loses its unstressed e before the ordinary vowel endings: hwæðres, hwæðrum, hwæðre, hwæðra. It retains e before the consonantal endings in hwæðerne, hwæðerre and hwæðerra. Its feminine nominative singular and neuter nominative/\allowbreak{}accusative plural are hwæðer. Do not copy the ōðru cells of ōðer into this paradigm. Attested alternatives include hwæðeres and feminine hwæðere.
% mdparagraph:P038:end

% mdparagraph:P039:begin
Hwæðer always distinguishes \textbf{two alternatives} in the use taught here. A singular noun can name one of two people or objects; plural morphology can distinguish two explicitly plural groups. It does not turn the choice into one among an unlimited number. The conjunction “whether” is a separate later use.
% mdparagraph:P039:end

% mdparagraph:P040:begin
In the second question below, two men are being compared.
% mdparagraph:P040:end

\reserve{8\baselineskip}
% example: NEW-U08-EXAMPLE-05
% mdparagraph:P041:begin
Hwilc mann hæfð ðæt scip? Hwæðer mann hæfð ðæt scip? Hūlic is ðæt scip?
% mdparagraph:P041:end

% mdparagraph:P042:begin
Which man has that ship? Which of the two men has that ship? What sort of ship is that?
% mdparagraph:P042:end

% mdparagraph:P043:begin
Hwilc and hwæðer modify nominative masculine singular mann. Hūlic is the strong neuter singular predicate referring to scip. The question word’s English meaning does not remove ordinary adjective agreement.
% mdparagraph:P043:end

\reserve{7\baselineskip}
% mdparagraph:P044:begin
\section*{3. Some, any, each, and such}
% mdparagraph:P044:end

% teaching: NEW-U08-GRAM-u08-indefinites
% mdparagraph:P045:begin
The following words also take complete strong adjective inflection. \textbf{Sum} means “a certain” in the singular and “some” in the plural. \textbf{Ǣnig} means “any” in the ordinary direct questions used here. \textbf{Gehwilc} means “each, every”; its ordinary use in this unit is with a singular noun. \textbf{Swilc} means “such”, referring to a kind already understood from the context.
% mdparagraph:P045:end

% teaching: NEW-U08-VOC-sum
% teaching: NEW-U08-PAR-SUM-STRONG
\reserve{7\baselineskip}
% mdparagraph:P046:begin
\subsection*{sum: strong declension, singular}
% mdparagraph:P046:end

% table: NEW-U08-PAR-SUM-STRONG-SG
% mdtable:T10:begin
\par\begingroup
\fontsize{10.5}{13.2}\selectfont
\setlength{\LTleft}{0pt}\setlength{\LTright}{\fill}
\noindent
\begin{tabular}{@{}>{\raggedright\arraybackslash}p{28.89mm}>{\raggedright\arraybackslash}p{28.89mm}>{\raggedright\arraybackslash}p{28.89mm}>{\raggedright\arraybackslash}p{28.89mm}@{}}
\toprule
\textbf{Case} & \textbf{Masculine} & \textbf{Neuter} & \textbf{Feminine} \\
\midrule
Nominative & sum & sum & sumu \\
Genitive & sumes & sumes & sumre \\
Dative & sumum & sumum & sumre \\
Accusative & sumne & sum & sume \\
\bottomrule
\end{tabular}
\par\endgroup

% mdtable:T10:end
\reserve{7\baselineskip}
% mdparagraph:P047:begin
\subsection*{sum: strong declension, plural}
% mdparagraph:P047:end

% table: NEW-U08-PAR-SUM-STRONG-PL
% mdtable:T11:begin
\par\begingroup
\fontsize{10.5}{13.2}\selectfont
\setlength{\LTleft}{0pt}\setlength{\LTright}{\fill}
\noindent
\begin{tabular}{@{}>{\raggedright\arraybackslash}p{28.89mm}>{\raggedright\arraybackslash}p{28.89mm}>{\raggedright\arraybackslash}p{28.89mm}>{\raggedright\arraybackslash}p{28.89mm}@{}}
\toprule
\textbf{Case} & \textbf{Masculine} & \textbf{Neuter} & \textbf{Feminine} \\
\midrule
Nominative & sume & sumu & suma \\
Genitive & sumra & sumra & sumra \\
Dative & sumum & sumum & sumum \\
Accusative & sume & sumu & suma \\
\bottomrule
\end{tabular}
\par\endgroup

% mdtable:T11:end
% mdparagraph:P048:begin
Short-stem sum has -u in feminine nominative singular and neuter nominative/\allowbreak{}accusative plural: sumu. The genitive and dative endings still mark their ordinary functions.
% mdparagraph:P048:end

\reserve{8\baselineskip}
% example: NEW-U08-EXAMPLE-06
% mdparagraph:P049:begin
Sum mann hæfð scip. Sume męnn habbað scipu.
% mdparagraph:P049:end

% mdparagraph:P050:begin
A certain man has a ship. Some men have ships.
% mdparagraph:P050:end

% mdparagraph:P051:begin
Sum mann is singular, with hæfð; sume męnn is plural, with habbað. The object likewise distinguishes scip from plural scipu. Sum supplies “a certain”; the earlier cardinal ān still means “one”.
% mdparagraph:P051:end

% teaching: NEW-U08-VOC-aenig
% teaching: NEW-U08-PAR-AENIG-STRONG
\reserve{7\baselineskip}
% mdparagraph:P052:begin
\subsection*{ǣnig: strong declension, singular}
% mdparagraph:P052:end

% table: NEW-U08-PAR-AENIG-STRONG-SG
% mdtable:T12:begin
\par\begingroup
\fontsize{10.5}{13.2}\selectfont
\setlength{\LTleft}{0pt}\setlength{\LTright}{\fill}
\noindent
\begin{tabular}{@{}>{\raggedright\arraybackslash}p{28.89mm}>{\raggedright\arraybackslash}p{28.89mm}>{\raggedright\arraybackslash}p{28.89mm}>{\raggedright\arraybackslash}p{28.89mm}@{}}
\toprule
\textbf{Case} & \textbf{Masculine} & \textbf{Neuter} & \textbf{Feminine} \\
\midrule
Nominative & ǣnig & ǣnig & ǣnigu \\
Genitive & ǣniges & ǣniges & ǣnigre \\
Dative & ǣnigum & ǣnigum & ǣnigre \\
Accusative & ǣnigne & ǣnig & ǣnige \\
\bottomrule
\end{tabular}
\par\endgroup

% mdtable:T12:end
\reserve{7\baselineskip}
% mdparagraph:P053:begin
\subsection*{ǣnig: strong declension, plural}
% mdparagraph:P053:end

% table: NEW-U08-PAR-AENIG-STRONG-PL
% mdtable:T13:begin
\par\begingroup
\fontsize{10.5}{13.2}\selectfont
\setlength{\LTleft}{0pt}\setlength{\LTright}{\fill}
\noindent
\begin{tabular}{@{}>{\raggedright\arraybackslash}p{28.89mm}>{\raggedright\arraybackslash}p{28.89mm}>{\raggedright\arraybackslash}p{28.89mm}>{\raggedright\arraybackslash}p{28.89mm}@{}}
\toprule
\textbf{Case} & \textbf{Masculine} & \textbf{Neuter} & \textbf{Feminine} \\
\midrule
Nominative & ǣnige & ǣnigu & ǣniga \\
Genitive & ǣnigra & ǣnigra & ǣnigra \\
Dative & ǣnigum & ǣnigum & ǣnigum \\
Accusative & ǣnige & ǣnigu & ǣniga \\
\bottomrule
\end{tabular}
\par\endgroup

% mdtable:T13:end
% mdparagraph:P054:begin
The selected ǣnig forms retain i: ǣniges, ǣnigum, ǣnigre, ǣnigra. The listed -u cells are ǣnigu; analogical ǣnig without -u is a documented alternative. Syncope in disyllabic adjectives varies, so do not mechanically contract every -ig adjective by analogy with another word. These questions teach a bounded use of “any”; they do not establish its distribution in every negative or conditional construction.
% mdparagraph:P054:end

% teaching: NEW-U08-VOC-gehwilc
% teaching: NEW-U08-PAR-GEHWILC-STRONG
\reserve{7\baselineskip}
% mdparagraph:P055:begin
\subsection*{gehwilc: strong declension, singular}
% mdparagraph:P055:end

% table: NEW-U08-PAR-GEHWILC-STRONG-SG
% mdtable:T14:begin
\par\begingroup
\fontsize{10.5}{13.2}\selectfont
\setlength{\LTleft}{0pt}\setlength{\LTright}{\fill}
\noindent
\begin{tabular}{@{}>{\raggedright\arraybackslash}p{28.89mm}>{\raggedright\arraybackslash}p{28.89mm}>{\raggedright\arraybackslash}p{28.89mm}>{\raggedright\arraybackslash}p{28.89mm}@{}}
\toprule
\textbf{Case} & \textbf{Masculine} & \textbf{Neuter} & \textbf{Feminine} \\
\midrule
Nominative & gehwilc & gehwilc & gehwilc \\
Genitive & gehwilces & gehwilces & gehwilcre \\
Dative & gehwilcum & gehwilcum & gehwilcre \\
Accusative & gehwilcne & gehwilc & gehwilce \\
\bottomrule
\end{tabular}
\par\endgroup

% mdtable:T14:end
\reserve{7\baselineskip}
% mdparagraph:P056:begin
\subsection*{gehwilc: strong declension, plural}
% mdparagraph:P056:end

% table: NEW-U08-PAR-GEHWILC-STRONG-PL
% mdtable:T15:begin
\par\begingroup
\fontsize{10.5}{13.2}\selectfont
\setlength{\LTleft}{0pt}\setlength{\LTright}{\fill}
\noindent
\begin{tabular}{@{}>{\raggedright\arraybackslash}p{28.89mm}>{\raggedright\arraybackslash}p{28.89mm}>{\raggedright\arraybackslash}p{28.89mm}>{\raggedright\arraybackslash}p{28.89mm}@{}}
\toprule
\textbf{Case} & \textbf{Masculine} & \textbf{Neuter} & \textbf{Feminine} \\
\midrule
Nominative & gehwilce & gehwilc & gehwilca \\
Genitive & gehwilcra & gehwilcra & gehwilcra \\
Dative & gehwilcum & gehwilcum & gehwilcum \\
Accusative & gehwilce & gehwilc & gehwilca \\
\bottomrule
\end{tabular}
\par\endgroup

% mdtable:T15:end
% mdparagraph:P057:begin
Gehwilc follows the same strong pattern as hwilc. A singular noun meaning “each man” or “each companion” takes a singular verb. Its complete plural morphology is given for form analysis, but the plural table does not make gehwilc a general substitute for eall “all” in every context. Partitive expressions such as “each of” come later.
% mdparagraph:P057:end

\reserve{8\baselineskip}
% example: NEW-U08-EXAMPLE-07
% mdparagraph:P058:begin
Hæfst ðū ǣnigne fisc? Gehwilc mann hæfð scip.
% mdparagraph:P058:end

% mdparagraph:P059:begin
Do you have any fish? Each man has a ship.
% mdparagraph:P059:end

% mdparagraph:P060:begin
Ǣnigne agrees with singular masculine accusative fisc; the English “any fish” here asks about a fish, not an Old English plural. Gehwilc mann is nominative singular.
% mdparagraph:P060:end

% teaching: NEW-U08-VOC-swilc
% teaching: NEW-U08-PAR-SWILC-STRONG
\reserve{7\baselineskip}
% mdparagraph:P061:begin
\subsection*{swilc: strong declension, singular}
% mdparagraph:P061:end

% table: NEW-U08-PAR-SWILC-STRONG-SG
% mdtable:T16:begin
\par\begingroup
\fontsize{10.5}{13.2}\selectfont
\setlength{\LTleft}{0pt}\setlength{\LTright}{\fill}
\noindent
\begin{tabular}{@{}>{\raggedright\arraybackslash}p{28.89mm}>{\raggedright\arraybackslash}p{28.89mm}>{\raggedright\arraybackslash}p{28.89mm}>{\raggedright\arraybackslash}p{28.89mm}@{}}
\toprule
\textbf{Case} & \textbf{Masculine} & \textbf{Neuter} & \textbf{Feminine} \\
\midrule
Nominative & swilc & swilc & swilc \\
Genitive & swilces & swilces & swilcre \\
Dative & swilcum & swilcum & swilcre \\
Accusative & swilcne & swilc & swilce \\
\bottomrule
\end{tabular}
\par\endgroup

% mdtable:T16:end
\reserve{7\baselineskip}
% mdparagraph:P062:begin
\subsection*{swilc: strong declension, plural}
% mdparagraph:P062:end

% table: NEW-U08-PAR-SWILC-STRONG-PL
% mdtable:T17:begin
\par\begingroup
\fontsize{10.5}{13.2}\selectfont
\setlength{\LTleft}{0pt}\setlength{\LTright}{\fill}
\noindent
\begin{tabular}{@{}>{\raggedright\arraybackslash}p{28.89mm}>{\raggedright\arraybackslash}p{28.89mm}>{\raggedright\arraybackslash}p{28.89mm}>{\raggedright\arraybackslash}p{28.89mm}@{}}
\toprule
\textbf{Case} & \textbf{Masculine} & \textbf{Neuter} & \textbf{Feminine} \\
\midrule
Nominative & swilce & swilc & swilca \\
Genitive & swilcra & swilcra & swilcra \\
Dative & swilcum & swilcum & swilcum \\
Accusative & swilce & swilc & swilca \\
\bottomrule
\end{tabular}
\par\endgroup

% mdtable:T17:end
% mdparagraph:P063:begin
Swilc keeps zero ending in feminine nominative singular and neuter nominative/\allowbreak{}accusative plural, as do hwilc, hūlic and gehwilc. In the selected indefinite use, a following ordinary descriptive adjective is strong. Its case, gender and number still agree with the same noun.
% mdparagraph:P063:end

\reserve{8\baselineskip}
% example: NEW-U08-EXAMPLE-10
% mdparagraph:P064:begin
Swilc mann is gōd. Swilce męnn sindon gōde.
% mdparagraph:P064:end

% mdparagraph:P065:begin
Such a man is good. Such men are good.
% mdparagraph:P065:end

% mdparagraph:P066:begin
Here swilce is an inflected masculine plural adjective modifying męnn. The earlier \textbf{swilce “likewise”} is an invariant adverb. Identical spelling does not make the two uses grammatically identical.
% mdparagraph:P066:end

\reserve{7\baselineskip}
% mdparagraph:P067:begin
\section*{4. Other, same, and self}
% mdparagraph:P067:end

% teaching: NEW-U08-VOC-other
% teaching: NEW-U08-GRAM-u08-other-sense
% mdparagraph:P068:begin
The earlier \textbf{ōðer “second”} also means \textbf{“other, another”}. Context determines the sense. Its strong forms are retained even after an article; this lexical pattern does not acquire a new weak declension.
% mdparagraph:P068:end

% teaching: NEW-U08-PAR-OTHER
\reserve{7\baselineskip}
% mdparagraph:P069:begin
\subsection*{ōðer: other /\allowbreak{} another, singular}
% mdparagraph:P069:end

% table: NEW-U08-PAR-OTHER-SG
% mdtable:T18:begin
\par\begingroup
\fontsize{10.5}{13.2}\selectfont
\setlength{\LTleft}{0pt}\setlength{\LTright}{\fill}
\noindent
\begin{tabular}{@{}>{\raggedright\arraybackslash}p{28.89mm}>{\raggedright\arraybackslash}p{28.89mm}>{\raggedright\arraybackslash}p{28.89mm}>{\raggedright\arraybackslash}p{28.89mm}@{}}
\toprule
\textbf{Case} & \textbf{Masculine} & \textbf{Neuter} & \textbf{Feminine} \\
\midrule
Nominative & ōðer & ōðer & ōðru \\
Genitive & ōðres & ōðres & ōðerre \\
Dative & ōðrum & ōðrum & ōðerre \\
Accusative & ōðerne & ōðer & ōðre \\
\bottomrule
\end{tabular}
\par\endgroup

% mdtable:T18:end
\reserve{7\baselineskip}
% mdparagraph:P070:begin
\subsection*{ōðer: other /\allowbreak{} another, plural}
% mdparagraph:P070:end

% table: NEW-U08-PAR-OTHER-PL
% mdtable:T19:begin
\par\begingroup
\fontsize{10.5}{13.2}\selectfont
\setlength{\LTleft}{0pt}\setlength{\LTright}{\fill}
\noindent
\begin{tabular}{@{}>{\raggedright\arraybackslash}p{28.89mm}>{\raggedright\arraybackslash}p{28.89mm}>{\raggedright\arraybackslash}p{28.89mm}>{\raggedright\arraybackslash}p{28.89mm}@{}}
\toprule
\textbf{Case} & \textbf{Masculine} & \textbf{Neuter} & \textbf{Feminine} \\
\midrule
Nominative & ōðre & ōðru & ōðra \\
Genitive & ōðerra & ōðerra & ōðerra \\
Dative & ōðrum & ōðrum & ōðrum \\
Accusative & ōðre & ōðru & ōðra \\
\bottomrule
\end{tabular}
\par\endgroup

% mdtable:T19:end
% mdparagraph:P071:begin
Retain ōðr- before the ordinary vowel endings, but ōðer- in ōðerne, ōðerre and ōðerra. In particular, feminine nominative singular and neuter nominative/\allowbreak{}accusative plural are ōðru. The complete earlier morphology is reused with the new sense.
% mdparagraph:P071:end

% teaching: NEW-U08-VOC-ilca
% teaching: NEW-U08-GRAM-u08-ilca
% mdparagraph:P072:begin
\textbf{Ilca “same” is weak only.} It ordinarily follows the article. Any following ordinary descriptive adjective is also weak in that definite phrase. Ilca has no strong declension.
% mdparagraph:P072:end

% teaching: NEW-U08-PAR-ILCA-WEAK
\reserve{7\baselineskip}
% mdparagraph:P073:begin
\subsection*{ilca: weak declension, singular}
% mdparagraph:P073:end

% table: NEW-U08-PAR-ILCA-WEAK-SG
% mdtable:T20:begin
\par\begingroup
\fontsize{10.5}{13.2}\selectfont
\setlength{\LTleft}{0pt}\setlength{\LTright}{\fill}
\noindent
\begin{tabular}{@{}>{\raggedright\arraybackslash}p{28.89mm}>{\raggedright\arraybackslash}p{28.89mm}>{\raggedright\arraybackslash}p{28.89mm}>{\raggedright\arraybackslash}p{28.89mm}@{}}
\toprule
\textbf{Case} & \textbf{Masculine} & \textbf{Neuter} & \textbf{Feminine} \\
\midrule
Nominative & ilca & ilce & ilce \\
Genitive & ilcan & ilcan & ilcan \\
Dative & ilcan & ilcan & ilcan \\
Accusative & ilcan & ilce & ilcan \\
\bottomrule
\end{tabular}
\par\endgroup

% mdtable:T20:end
\reserve{7\baselineskip}
% mdparagraph:P074:begin
\subsection*{ilca: weak declension, plural}
% mdparagraph:P074:end

% table: NEW-U08-PAR-ILCA-WEAK-PL
% mdtable:T21:begin
\par\begingroup
\fontsize{10.5}{13.2}\selectfont
\setlength{\LTleft}{0pt}\setlength{\LTright}{\fill}
\noindent
\begin{tabular}{@{}>{\raggedright\arraybackslash}p{28.89mm}>{\raggedright\arraybackslash}p{28.89mm}>{\raggedright\arraybackslash}p{28.89mm}>{\raggedright\arraybackslash}p{28.89mm}@{}}
\toprule
\textbf{Case} & \textbf{Masculine} & \textbf{Neuter} & \textbf{Feminine} \\
\midrule
Nominative & ilcan & ilcan & ilcan \\
Genitive & ilcra & ilcra & ilcra \\
Dative & ilcum & ilcum & ilcum \\
Accusative & ilcan & ilcan & ilcan \\
\bottomrule
\end{tabular}
\par\endgroup

% mdtable:T21:end
% mdparagraph:P075:begin
The stem ilc- takes masculine nominative singular -a; feminine/\allowbreak{}neuter nominative singular and neuter accusative singular -e; the remaining singular endings are -an. Plural nominative and accusative have -an, genitive -ra and dative -um.
% mdparagraph:P075:end

\reserve{8\baselineskip}
% example: NEW-U08-EXAMPLE-08
% mdparagraph:P076:begin
Se ōðer mann hæfð ōðer scip. Se ilca mann hæfð ðæt scip.
% mdparagraph:P076:end

% mdparagraph:P077:begin
The other man has another ship. The same man has that ship.
% mdparagraph:P077:end

% mdparagraph:P078:begin
The first sentence distinguishes another man and another ship; it need not place them second in a numbered series. The second identifies the same man again. Both ōðer and ilca agree with the masculine nominative singular mann, but each follows its own declensional pattern.
% mdparagraph:P078:end

% teaching: NEW-U08-VOC-self
% teaching: NEW-U08-GRAM-u08-self
% mdparagraph:P079:begin
\textbf{Self} emphasizes or reinforces the person or thing referred to. It agrees with that referent in case, number and gender; it is not an invariant ending mechanically attached to an English-style pronoun compound. Both strong and weak forms occur. Strong self is used in the ordinary personal emphasis illustrated below; the weak system is retained separately, with the familiar weak endings.
% mdparagraph:P079:end

% teaching: NEW-U08-PAR-SELF-STRONG
\reserve{7\baselineskip}
% mdparagraph:P080:begin
\subsection*{self: strong declension, singular}
% mdparagraph:P080:end

% table: NEW-U08-PAR-SELF-STRONG-SG
% mdtable:T22:begin
\par\begingroup
\fontsize{10.5}{13.2}\selectfont
\setlength{\LTleft}{0pt}\setlength{\LTright}{\fill}
\noindent
\begin{tabular}{@{}>{\raggedright\arraybackslash}p{28.89mm}>{\raggedright\arraybackslash}p{28.89mm}>{\raggedright\arraybackslash}p{28.89mm}>{\raggedright\arraybackslash}p{28.89mm}@{}}
\toprule
\textbf{Case} & \textbf{Masculine} & \textbf{Neuter} & \textbf{Feminine} \\
\midrule
Nominative & self & self & self \\
Genitive & selfes & selfes & selfre \\
Dative & selfum & selfum & selfre \\
Accusative & selfne & self & selfe \\
\bottomrule
\end{tabular}
\par\endgroup

% mdtable:T22:end
\reserve{7\baselineskip}
% mdparagraph:P081:begin
\subsection*{self: strong declension, plural}
% mdparagraph:P081:end

% table: NEW-U08-PAR-SELF-STRONG-PL
% mdtable:T23:begin
\par\begingroup
\fontsize{10.5}{13.2}\selectfont
\setlength{\LTleft}{0pt}\setlength{\LTright}{\fill}
\noindent
\begin{tabular}{@{}>{\raggedright\arraybackslash}p{28.89mm}>{\raggedright\arraybackslash}p{28.89mm}>{\raggedright\arraybackslash}p{28.89mm}>{\raggedright\arraybackslash}p{28.89mm}@{}}
\toprule
\textbf{Case} & \textbf{Masculine} & \textbf{Neuter} & \textbf{Feminine} \\
\midrule
Nominative & selfe & self & selfa \\
Genitive & selfra & selfra & selfra \\
Dative & selfum & selfum & selfum \\
Accusative & selfe & self & selfa \\
\bottomrule
\end{tabular}
\par\endgroup

% mdtable:T23:end
% teaching: NEW-U08-PAR-SELF-WEAK
\reserve{7\baselineskip}
% mdparagraph:P082:begin
\subsection*{self: weak declension, singular}
% mdparagraph:P082:end

% table: NEW-U08-PAR-SELF-WEAK-SG
% mdtable:T24:begin
\par\begingroup
\fontsize{10.5}{13.2}\selectfont
\setlength{\LTleft}{0pt}\setlength{\LTright}{\fill}
\noindent
\begin{tabular}{@{}>{\raggedright\arraybackslash}p{28.89mm}>{\raggedright\arraybackslash}p{28.89mm}>{\raggedright\arraybackslash}p{28.89mm}>{\raggedright\arraybackslash}p{28.89mm}@{}}
\toprule
\textbf{Case} & \textbf{Masculine} & \textbf{Neuter} & \textbf{Feminine} \\
\midrule
Nominative & selfa & selfe & selfe \\
Genitive & selfan & selfan & selfan \\
Dative & selfan & selfan & selfan \\
Accusative & selfan & selfe & selfan \\
\bottomrule
\end{tabular}
\par\endgroup

% mdtable:T24:end
\reserve{7\baselineskip}
% mdparagraph:P083:begin
\subsection*{self: weak declension, plural}
% mdparagraph:P083:end

% table: NEW-U08-PAR-SELF-WEAK-PL
% mdtable:T25:begin
\par\begingroup
\fontsize{10.5}{13.2}\selectfont
\setlength{\LTleft}{0pt}\setlength{\LTright}{\fill}
\noindent
\begin{tabular}{@{}>{\raggedright\arraybackslash}p{28.89mm}>{\raggedright\arraybackslash}p{28.89mm}>{\raggedright\arraybackslash}p{28.89mm}>{\raggedright\arraybackslash}p{28.89mm}@{}}
\toprule
\textbf{Case} & \textbf{Masculine} & \textbf{Neuter} & \textbf{Feminine} \\
\midrule
Nominative & selfan & selfan & selfan \\
Genitive & selfra & selfra & selfra \\
Dative & selfum & selfum & selfum \\
Accusative & selfan & selfan & selfan \\
\bottomrule
\end{tabular}
\par\endgroup

% mdtable:T25:end
\reserve{8\baselineskip}
% example: NEW-U08-EXAMPLE-09
% mdparagraph:P084:begin
Hē self geseah ðæt scip.
% mdparagraph:P084:end

% mdparagraph:P085:begin
He himself saw that ship.
% mdparagraph:P085:end

% mdparagraph:P086:begin
Self is strong nominative masculine singular, emphasizing the subject hē. A personal pronoun can already have reflexive reference without self. Adding self reinforces it: an accusative reflexive takes accusative agreement, and a dative reflexive takes dative agreement. For a male accusative referent the strong form is selfne; the first-person form mē can therefore be reinforced by selfne when the speaker is male. Do not choose self’s case from the English translation “myself” alone.
% mdparagraph:P086:end

\reserve{7\baselineskip}
% mdparagraph:P087:begin
\subsection*{Reference forms and limits}
% mdparagraph:P087:end

% teaching: NEW-U08-VOC-gehwa
% teaching: NEW-U08-GRAM-u08-gehwa-reference
% mdparagraph:P088:begin
\textbf{Gehwā /\allowbreak{} gehwæt} is an indefinite ge-compound of the interrogative pattern: “each one, everyone /\allowbreak{} everything”. Recognize the following complete singular paradigm. Its independent indefinite, partitive and distributive uses are reserved for later work; no plural is supplied.
% mdparagraph:P088:end

% teaching: NEW-U08-REF-GEHWA
\reserve{7\baselineskip}
% mdparagraph:P089:begin
\subsection*{gehwā /\allowbreak{} gehwæt: reference forms}
% mdparagraph:P089:end

% table: NEW-U08-REF-GEHWA
% mdtable:T26:begin
\par\begingroup
\fontsize{10.5}{13.2}\selectfont
\setlength{\LTleft}{0pt}\setlength{\LTright}{\fill}
\noindent
\begin{tabular}{@{}>{\raggedright\arraybackslash}p{25.54mm}>{\raggedright\arraybackslash}p{64.36mm}>{\raggedright\arraybackslash}p{28.48mm}@{}}
\toprule
\textbf{Case} & \textbf{Common masculine/\allowbreak{}feminine singular} & \textbf{Neuter singular} \\
\midrule
Nominative & gehwā & gehwæt \\
Genitive & gehwæs & gehwæs \\
Dative & gehwǣm & gehwǣm \\
Accusative & gehwone & gehwæt \\
\bottomrule
\end{tabular}
\par\endgroup

% mdtable:T26:end
% mdparagraph:P090:begin
The distinct strong instrumental forms below are also for recognition. Weak instrumental forms coincide with dative forms. These reference tables do not introduce productive instrumental constructions.
% mdparagraph:P090:end

% teaching: NEW-U08-REF-HWILC-INSTR
% teaching: NEW-U08-REF-HULIC-INSTR
% teaching: NEW-U08-REF-HWAETHER-INSTR
% teaching: NEW-U08-REF-SUM-INSTR
% teaching: NEW-U08-REF-AENIG-INSTR
% teaching: NEW-U08-REF-GEHWILC-INSTR
% teaching: NEW-U08-REF-SWILC-INSTR
% teaching: NEW-U08-REF-SELF-INSTR
% teaching: NEW-U08-REF-OTHER-INSTR
\reserve{7\baselineskip}
% mdparagraph:P091:begin
\subsection*{Instrumental reference forms}
% mdparagraph:P091:end

% table: NEW-U08-MATRIX-INSTRUMENTAL-REFERENCE
% mdtable:T27:begin
\par\begingroup
\fontsize{10.5}{13.2}\selectfont
\setlength{\LTleft}{0pt}\setlength{\LTright}{\fill}
\noindent
\begin{tabular}{@{}>{\raggedright\arraybackslash}p{34.61mm}>{\raggedright\arraybackslash}p{45.17mm}>{\raggedright\arraybackslash}p{38.59mm}@{}}
\toprule
\textbf{Word} & \textbf{Masculine singular} & \textbf{Neuter singular} \\
\midrule
hwilc & hwilce & hwilce \\
hūlic & hūlice & hūlice \\
hwæðer & hwæðre & hwæðre \\
sum & sume & sume \\
ǣnig & ǣnige & ǣnige \\
gehwilc & gehwilce & gehwilce \\
swilc & swilce & swilce \\
self & selfe & selfe \\
ōðer & ōðre & ōðre \\
\bottomrule
\end{tabular}
\par\endgroup

% mdtable:T27:end
% mdparagraph:P092:begin
The ordinary uses above leave several related systems for Unit 15: indefinite ān “a certain”, negative pronouns such as nān and nǣnig, either/\allowbreak{}neither families, indeclinable “anything/\allowbreak{}nothing” words, -hwega compounds, indefinite hwā “anyone”, and swā-hwā/\allowbreak{}swā-hwilc free-choice clauses. The impersonal use of man “one” is also separate from the count noun mann already learned. The earlier meanings and limits of eall, ǣlc and the possessives remain in force; sīn is not added here.
% mdparagraph:P092:end

\reserve{7\baselineskip}
% mdparagraph:P093:begin
\section*{5. Relative clauses}
% mdparagraph:P093:end

% teaching: NEW-U08-VOC-the-relative
% teaching: NEW-U08-VOC-se
% teaching: NEW-U08-GRAM-u08-relative-the
% teaching: NEW-U08-GRAM-u08-relative-demonstrative
% teaching: NEW-U08-GRAM-u08-relative-internal-case
% mdparagraph:P094:begin
A \textbf{relative clause} identifies or describes an antecedent in another clause. Old English can introduce it with invariant \textbf{ðe}, with an inflected member of \textbf{se/\allowbreak{}sēo/\allowbreak{}ðæt} alone, or with that inflected member followed by \textbf{ðe}. English often uses “who”, “which” or “that”; an English personal relative does not require the Old English interrogative hwā.
% mdparagraph:P094:end

\reserve{7\baselineskip}
% mdparagraph:P095:begin
\subsection*{Invariant relative particle}
% mdparagraph:P095:end

% table: NEW-U08-MATRIX-RELATIVE-PARTICLE
% mdtable:T28:begin
\par\begingroup
\fontsize{10.5}{13.2}\selectfont
\setlength{\LTleft}{0pt}\setlength{\LTright}{\fill}
\noindent
\begin{tabular}{@{}>{\raggedright\arraybackslash}p{92.06mm}>{\raggedright\arraybackslash}p{29.13mm}@{}}
\toprule
\textbf{Function} & \textbf{Form} \\
\midrule
Who /\allowbreak{} which /\allowbreak{} that, with role determined inside the clause & ðe \\
\bottomrule
\end{tabular}
\par\endgroup

% mdtable:T28:end
% mdparagraph:P096:begin
Ðe does not change its spelling for gender, number or case. Determine its role from the relative clause itself.
% mdparagraph:P096:end

\reserve{8\baselineskip}
% example: NEW-U08-EXAMPLE-11
% mdparagraph:P097:begin
Se mann, ðe ðæt scip hæfð, is mīn gefēra.
% mdparagraph:P097:end

% mdparagraph:P098:begin
The man who has that ship is my companion.
% mdparagraph:P098:end

% mdparagraph:P099:begin
Ðe is the subject of hæfð; the ship is its accusative object. The antecedent se mann is the subject of the main clause is mīn gefēra.
% mdparagraph:P099:end

\reserve{8\baselineskip}
% example: NEW-U08-EXAMPLE-12
% mdparagraph:P100:begin
Ic geseah ðæt scip, ðe se mann hæfð.
% mdparagraph:P100:end

% mdparagraph:P101:begin
I saw the ship that the man has.
% mdparagraph:P101:end

% mdparagraph:P102:begin
Here ðe is the object of hæfð, while se mann is the relative clause’s own nominative subject. The antecedent ðæt scip is the object of the main verb geseah. Invariant spelling does not erase these different functions.
% mdparagraph:P102:end

% mdparagraph:P103:begin
The inflected relative reuses the complete demonstrative forms. Its \textbf{gender and number} normally follow the antecedent. In the ordinary pattern used for composition here, its \textbf{case} follows its function inside the relative clause. The antecedent has the case required by its separate function in the main clause.
% mdparagraph:P103:end

% teaching: NEW-U08-PAR-RELATIVE-SE
\reserve{7\baselineskip}
% mdparagraph:P104:begin
\subsection*{se /\allowbreak{} sēo /\allowbreak{} ðæt: relative forms, singular}
% mdparagraph:P104:end

% table: NEW-U08-PAR-RELATIVE-SE-SG
% mdtable:T29:begin
\par\begingroup
\fontsize{10.5}{13.2}\selectfont
\setlength{\LTleft}{0pt}\setlength{\LTright}{\fill}
\noindent
\begin{tabular}{@{}>{\raggedright\arraybackslash}p{28.89mm}>{\raggedright\arraybackslash}p{28.89mm}>{\raggedright\arraybackslash}p{28.89mm}>{\raggedright\arraybackslash}p{28.89mm}@{}}
\toprule
\textbf{Case} & \textbf{Masculine} & \textbf{Neuter} & \textbf{Feminine} \\
\midrule
Nominative & se & ðæt & sēo \\
Genitive & ðæs & ðæs & ðǣre \\
Dative & ðǣm & ðǣm & ðǣre \\
Accusative & ðone & ðæt & ðā \\
\bottomrule
\end{tabular}
\par\endgroup

% mdtable:T29:end
\reserve{7\baselineskip}
% mdparagraph:P105:begin
\subsection*{se /\allowbreak{} sēo /\allowbreak{} ðæt: relative forms, plural}
% mdparagraph:P105:end

% table: NEW-U08-PAR-RELATIVE-SE-PL
% mdtable:T30:begin
\par\begingroup
\fontsize{10.5}{13.2}\selectfont
\setlength{\LTleft}{0pt}\setlength{\LTright}{\fill}
\noindent
\begin{tabular}{@{}>{\raggedright\arraybackslash}p{28.89mm}>{\raggedright\arraybackslash}p{28.89mm}>{\raggedright\arraybackslash}p{28.89mm}>{\raggedright\arraybackslash}p{28.89mm}@{}}
\toprule
\textbf{Case} & \textbf{Masculine} & \textbf{Neuter} & \textbf{Feminine} \\
\midrule
Nominative & ðā & ðā & ðā \\
Genitive & ðāra & ðāra & ðāra \\
Dative & ðǣm & ðǣm & ðǣm \\
Accusative & ðā & ðā & ðā \\
\bottomrule
\end{tabular}
\par\endgroup

% mdtable:T30:end
\reserve{8\baselineskip}
% example: NEW-U08-EXAMPLE-13
% mdparagraph:P106:begin
Ic sęlle ðone fisc ðǣm męnn, se ðe mē healp.
% mdparagraph:P106:end

% mdparagraph:P107:begin
I give the fish to the man who helped me.
% mdparagraph:P107:end

% mdparagraph:P108:begin
Ðǣm męnn is dative, the recipient of sęlle. The relative se is nominative, the subject of healp. The person helped is dative mē. The antecedent’s dative does not make the relative subject dative.
% mdparagraph:P108:end

\reserve{8\baselineskip}
% example: NEW-U08-EXAMPLE-14
% mdparagraph:P109:begin
Ic geseah ðone mann, ðǣm ðe ic healp.
% mdparagraph:P109:end

% mdparagraph:P110:begin
I saw the man whom I helped.
% mdparagraph:P110:end

% mdparagraph:P111:begin
Ðone mann is accusative after geseah. Within the relative, ic is subject and helpan governs dative ðǣm. The relative’s gender and number still refer to the masculine singular man.
% mdparagraph:P111:end

\reserve{8\baselineskip}
% example: NEW-U08-EXAMPLE-16
% mdparagraph:P112:begin
Se mann, ðæs scip ic geseah, is gōd.
% mdparagraph:P112:end

% mdparagraph:P113:begin
The man whose ship I saw is good.
% mdparagraph:P113:end

% mdparagraph:P114:begin
Ðæs is genitive, expressing possession inside the relative: \textbf{whose ship}. Its antecedent mann is masculine singular and nominative in the main clause. Scip is neuter singular but accusative, the object of geseah; ic is the subject. The possessor’s antecedent determines the relative’s gender and number, not the possessed ship. A possessive relative therefore does not force its possessed noun into the genitive.
% mdparagraph:P114:end

% teaching: NEW-U08-GRAM-u08-relative-person-agreement
% mdparagraph:P115:begin
When a relative subject refers to a personal-pronoun antecedent, its verb can retain that antecedent’s person. In the selected second-person clauses, use the second-person verb form; the English word “who” does not automatically make the Old English verb third person.
% mdparagraph:P115:end

\reserve{8\baselineskip}
% example: NEW-U08-EXAMPLE-15
% mdparagraph:P116:begin
Ðū, ðe mē hulpe, sęle mē ðone fisc.
% mdparagraph:P116:end

% mdparagraph:P117:begin
You, who helped me, give me that fish.
% mdparagraph:P117:end

% mdparagraph:P118:begin
Hulpe is second-person singular preterite, agreeing with the antecedent ðū. Sęle is the singular command addressed to that person. With a plural antecedent the relative verb and the command must both change where their forms require it.
% mdparagraph:P118:end

% teaching: NEW-U08-GRAM-u08-relative-case-attraction
% teaching: NEW-U08-GRAM-u08-relative-personal-reference
% mdparagraph:P119:begin
\textbf{Recognition notes.} An inflected relative sometimes takes its antecedent’s case instead of the case expected from its own role. This is \textbf{case attraction}, an attested exception; it does not replace the ordinary internal-case pattern required in this unit’s compositions. A less common strategy uses ðe together with a personal pronoun whose inflection marks the internal function. Recognize that strategy when explained in a text; independent composition with it comes later. Neither note changes the earlier limits on personal genitives or instrumental syntax.
% mdparagraph:P119:end

\reserve{7\baselineskip}
% mdparagraph:P120:begin
\section*{6. Additional verb systems and lexical constructions}
% mdparagraph:P120:end

% teaching: NEW-U08-VOC-aeteowian
% teaching: NEW-U08-GRAM-u08-aeteowian-finite
% teaching: NEW-U08-GRAM-u08-appear-intransitive
\reserve{7\baselineskip}
% mdparagraph:P121:begin
\subsection*{Ætēowian: appear}
% mdparagraph:P121:end

% mdparagraph:P122:begin
\textbf{Ætēowian} is weak class II and \textbf{intransitive}: the thing appearing is the nominative subject. It is distinct from transitive ætīewan “reveal, display”. The complete present and preterite indicative and subjunctive are given below.
% mdparagraph:P122:end

% mdparagraph:P123:begin
Keep the lexical stem ætēow-. The present indicative has first-singular -ie, second-singular -ast, third-singular -að and plural -iað. The weak preterite has -ode, -odest, -ode and plural -odon. For the subjunctive, add singular -e or plural -en to the present stem ætēowi- or the preterite stem ætēowod-, giving the exact spellings shown.
% mdparagraph:P123:end

% teaching: NEW-U08-PAR-AETEOWIAN-FINITE
\reserve{7\baselineskip}
% mdparagraph:P124:begin
\subsection*{ætēowian: complete indicative}
% mdparagraph:P124:end

% table: NEW-U08-PAR-AETEOWIAN-FINITE-INDICATIVE
% mdtable:T31:begin
\par\begingroup
\fontsize{10.5}{13.2}\selectfont
\setlength{\LTleft}{0pt}\setlength{\LTright}{\fill}
\noindent
\begin{tabular}{@{}>{\raggedright\arraybackslash}p{35.53mm}>{\raggedright\arraybackslash}p{40.1mm}>{\raggedright\arraybackslash}p{42.75mm}@{}}
\toprule
\textbf{Person and number} & \textbf{Present indicative} & \textbf{Preterite indicative} \\
\midrule
1 singular & ætēowie & ætēowode \\
2 singular & ætēowast & ætēowodest \\
3 singular & ætēowað & ætēowode \\
1 plural & ætēowiað & ætēowodon \\
2 plural & ætēowiað & ætēowodon \\
3 plural & ætēowiað & ætēowodon \\
\bottomrule
\end{tabular}
\par\endgroup

% mdtable:T31:end
\reserve{7\baselineskip}
% mdparagraph:P125:begin
\subsection*{ætēowian: complete subjunctive}
% mdparagraph:P125:end

% table: NEW-U08-PAR-AETEOWIAN-FINITE-SUBJUNCTIVE
% mdtable:T32:begin
\par\begingroup
\fontsize{10.5}{13.2}\selectfont
\setlength{\LTleft}{0pt}\setlength{\LTright}{\fill}
\noindent
\begin{tabular}{@{}>{\raggedright\arraybackslash}p{42.16mm}>{\raggedright\arraybackslash}p{36.97mm}>{\raggedright\arraybackslash}p{39.25mm}@{}}
\toprule
\textbf{Number (all three persons)} & \textbf{Present subjunctive} & \textbf{Preterite subjunctive} \\
\midrule
Singular: persons 1, 2, 3 & ætēowie & ætēowode \\
Plural: persons 1, 2, 3 & ætēowien & ætēowoden \\
\bottomrule
\end{tabular}
\par\endgroup

% mdtable:T32:end
% mdparagraph:P126:begin
The present subjunctive can express “let … appear” in the wish or direction already learned. No separate lǣtan is needed for that meaning. Cook’s actual form in the Creation text is \textbf{ætēowie}, without g. The documented weak-II i/\allowbreak{}ig(e)/\allowbreak{}g variation does not replace the selected i spellings in these tables.
% mdparagraph:P126:end

\reserve{7\baselineskip}
% mdparagraph:P127:begin
\subsection*{ætēowian reference citations}
% mdparagraph:P127:end

% table: NEW-U08-MATRIX-AETEOWIAN-CITATIONS
% mdtable:T33:begin
\par\begingroup
\fontsize{10.5}{13.2}\selectfont
\setlength{\LTleft}{0pt}\setlength{\LTright}{\fill}
\noindent
\begin{tabular}{@{}>{\raggedright\arraybackslash}p{79.47mm}>{\raggedright\arraybackslash}p{41.72mm}@{}}
\toprule
\textbf{Function} & \textbf{Form} \\
\midrule
Dictionary infinitive & ætēowian \\
Past-participle citation & ætēowod \\
\bottomrule
\end{tabular}
\par\endgroup

% mdtable:T33:end
% mdparagraph:P128:begin
The infinitive and past participle in this small table are dictionary reference forms. This unit’s new productive system is finite indicative and subjunctive; it does not add commands or new nonfinite constructions of ætēowian.
% mdparagraph:P128:end

\reserve{7\baselineskip}
% mdparagraph:P129:begin
\subsection*{Scīnan: subjunctive}
% mdparagraph:P129:end

% teaching: NEW-U08-VOC-scinan
% teaching: NEW-U08-GRAM-u08-scinan-subj
% mdparagraph:P130:begin
The indicative of strong-I \textbf{scīnan “shine”} is familiar. Its complete subjunctive uses present scīn- and the preterite plural stem scin-. The long ī belongs to the present; the past subjunctive has short i, not the scān- of the preterite indicative singular.
% mdparagraph:P130:end

% teaching: NEW-U08-PAR-SCINAN-SUBJ
\reserve{7\baselineskip}
% mdparagraph:P131:begin
\subsection*{scīnan: subjunctive}
% mdparagraph:P131:end

% table: NEW-U08-PAR-SCINAN-SUBJ
% mdtable:T34:begin
\par\begingroup
\fontsize{10.5}{13.2}\selectfont
\setlength{\LTleft}{0pt}\setlength{\LTright}{\fill}
\noindent
\begin{tabular}{@{}>{\raggedright\arraybackslash}p{42.16mm}>{\raggedright\arraybackslash}p{36.97mm}>{\raggedright\arraybackslash}p{39.25mm}@{}}
\toprule
\textbf{Number (all three persons)} & \textbf{Present subjunctive} & \textbf{Preterite subjunctive} \\
\midrule
Singular: persons 1, 2, 3 & scīne & scine \\
Plural: persons 1, 2, 3 & scīnen & scinen \\
\bottomrule
\end{tabular}
\par\endgroup

% mdtable:T34:end
% mdparagraph:P132:begin
Each singular or plural row covers all three persons, in both tenses. The subject and the clause determine the person; the subjunctive ending itself distinguishes the two numbers.
% mdparagraph:P132:end

\reserve{7\baselineskip}
% mdparagraph:P133:begin
\subsection*{Geblētsian: commands}
% mdparagraph:P133:end

% teaching: NEW-U08-VOC-gebletsian
% teaching: NEW-U08-GRAM-u08-gebletsian-imperative
% mdparagraph:P134:begin
Weak-II \textbf{geblētsian “bless”} has singular imperative -a and plural -iað. Retain the prefix ge-. The number of people being blessed is the object’s number; the command form depends on the number of people addressed.
% mdparagraph:P134:end

% teaching: NEW-U08-PAR-GEBLETSIAN-IMP
\reserve{7\baselineskip}
% mdparagraph:P135:begin
\subsection*{geblētsian: commands}
% mdparagraph:P135:end

% table: NEW-U08-PAR-GEBLETSIAN-IMP
% mdtable:T35:begin
\par\begingroup
\fontsize{10.5}{13.2}\selectfont
\setlength{\LTleft}{0pt}\setlength{\LTright}{\fill}
\noindent
\begin{tabular}{@{}>{\raggedright\arraybackslash}p{77.12mm}>{\raggedright\arraybackslash}p{44.07mm}@{}}
\toprule
\textbf{Second-person addressee} & \textbf{Imperative} \\
\midrule
One person & geblētsa \\
Two or more people & geblētsiað \\
\bottomrule
\end{tabular}
\par\endgroup

% mdtable:T35:end
\reserve{7\baselineskip}
% mdparagraph:P136:begin
\subsection*{Līf and the cognate accusative}
% mdparagraph:P136:end

% teaching: NEW-U08-VOC-lif
% teaching: NEW-U08-GRAM-u08-lif-plural
% teaching: NEW-U08-GRAM-u08-cognate-accusative
% mdparagraph:P137:begin
Neuter \textbf{līf “life”} can name an individual life, with the ordinary plural “lives”. Its long stem has no ending in nominative or accusative plural. The full declension is repeated here for this use.
% mdparagraph:P137:end

% teaching: NEW-U08-PAR-LIF
\reserve{7\baselineskip}
% mdparagraph:P138:begin
\subsection*{līf: complete declension}
% mdparagraph:P138:end

% table: NEW-U08-PAR-LIF
% mdtable:T36:begin
\par\begingroup
\fontsize{10.5}{13.2}\selectfont
\setlength{\LTleft}{0pt}\setlength{\LTright}{\fill}
\noindent
\begin{tabular}{@{}>{\raggedright\arraybackslash}p{39.46mm}>{\raggedright\arraybackslash}p{39.46mm}>{\raggedright\arraybackslash}p{39.46mm}@{}}
\toprule
\textbf{Case} & \textbf{Singular} & \textbf{Plural} \\
\midrule
Nominative & līf & līf \\
Genitive & līfes & līfa \\
Dative & līfe & līfum \\
Accusative & līf & līf \\
\bottomrule
\end{tabular}
\par\endgroup

% mdtable:T36:end
% mdparagraph:P139:begin
\textbf{Libban} can take the related noun līf as a \textbf{cognate accusative}: “live one’s life”, “live our lives”. The object names the life being lived; it is not a locative phrase. In “our lives”, the possessive is neuter accusative plural \textbf{ūru}, agreeing with plural līf. This is a particular construction of libban with līf, not a rule permitting an arbitrary object after every intransitive verb. Cook I.8.8 refers to his §168.1.
% mdparagraph:P139:end

\reserve{7\baselineskip}
% mdparagraph:P140:begin
\subsection*{Fill; have in one’s control}
% mdparagraph:P140:end

% teaching: NEW-U08-VOC-gefyllan
% teaching: NEW-U08-GRAM-u08-gefyllan-fill
% mdparagraph:P141:begin
\textbf{Gefyllan}, weak I, now also has the ordinary active sense \textbf{“fill”}, with the thing or place filled in the accusative. Keep this sense distinct from the earlier “finish, accomplish”, as in finishing a work. The already learned finite and imperative forms serve both senses; no new endings are required.
% mdparagraph:P141:end

% teaching: NEW-U08-VOC-on
% teaching: NEW-U08-GRAM-u08-control-on
% mdparagraph:P142:begin
\textbf{Habban + accusative controlled entity + on + dative possessive/\allowbreak{}gewealde} means “have in one’s control”. Geweald is neuter; its dative singular is \textbf{gewealde}. An agreeing possessive takes dative neuter singular: for one addressee, \textbf{ðīnum}. The invariant possessor \textbf{his} keeps its form. This is the established control expression, not a spatial translation of English “under”. The controlled persons or things remain the accusative object of habban.
% mdparagraph:P142:end

\reserve{7\baselineskip}
% mdparagraph:P143:begin
\section*{7. Practice with forms and clauses}
% mdparagraph:P143:end

% mdparagraph:P144:begin
Use the selected table forms when a task requests a paradigm. In composition, preserve the stated person, number, tense and construction. The questions and relatives use the ordinary case patterns explained above.
% mdparagraph:P144:end

\reserve{7\baselineskip}
% mdparagraph:P145:begin
\subsection*{Practice A. Complete forms}
% mdparagraph:P145:end

\reserve{5\baselineskip}
\par\begin{minipage}{\linewidth}\raggedright
% mdparagraph:P146:begin
\textbf{A1.}
% mdparagraph:P146:end

% item: NEW-U08-A01
% mdparagraph:P147:begin
Give the nominative, genitive, dative and accusative singular of hwā and hwæt. Label each form and distinguish the common personal column from the neuter column.
% mdparagraph:P147:end

\end{minipage}\par\addvspace{4pt}

\par\begin{minipage}{\linewidth}\raggedright
% mdparagraph:P148:begin
\textbf{A2.}
% mdparagraph:P148:end

% item: NEW-U08-A02
% mdparagraph:P149:begin
Give the strong singular genitive, dative and accusative of hwilc in all three genders, and its nominative plural in all three genders.
% mdparagraph:P149:end

\end{minipage}\par\addvspace{4pt}

\par\begin{minipage}{\linewidth}\raggedright
% mdparagraph:P150:begin
\textbf{A3.}
% mdparagraph:P150:end

% item: NEW-U08-A03
% mdparagraph:P151:begin
Give the strong masculine accusative singular, feminine dative singular, neuter nominative plural and genitive plural of hwæðer and hūlic. Label every form.
% mdparagraph:P151:end

\end{minipage}\par\addvspace{4pt}

\par\begin{minipage}{\linewidth}\raggedright
% mdparagraph:P152:begin
\textbf{A4.}
% mdparagraph:P152:end

% item: NEW-U08-A04
% mdparagraph:P153:begin
Give the complete masculine singular and neuter plural declension of ilca. Explain why a strong alternative is not the model assigned here.
% mdparagraph:P153:end

\end{minipage}\par\addvspace{4pt}

\par\begin{minipage}{\linewidth}\raggedright
% mdparagraph:P154:begin
\textbf{A5.}
% mdparagraph:P154:end

% item: NEW-U08-A05
% mdparagraph:P155:begin
Give the nominative singular of self in all three genders, first strong and then weak. Also give its strong masculine accusative singular and strong dative plural.
% mdparagraph:P155:end

\end{minipage}\par\addvspace{4pt}

\par\begin{minipage}{\linewidth}\raggedright
% mdparagraph:P156:begin
\textbf{A6.}
% mdparagraph:P156:end

% item: NEW-U08-A06
% mdparagraph:P157:begin
Give the strong masculine accusative singular, feminine genitive singular, neuter nominative plural and dative plural of sum, ǣnig, gehwilc and swilc.
% mdparagraph:P157:end

\end{minipage}\par\addvspace{4pt}

\par\begin{minipage}{\linewidth}\raggedright
% mdparagraph:P158:begin
\textbf{A7.}
% mdparagraph:P158:end

% item: NEW-U08-A07
% mdparagraph:P159:begin
Give the complete singular and plural declension of līf. Explain why the plural in “our lives” is an ordinary contextual plural in this unit.
% mdparagraph:P159:end

\end{minipage}\par\addvspace{4pt}

\par\begin{minipage}{\linewidth}\raggedright
% mdparagraph:P160:begin
\textbf{A8.}
% mdparagraph:P160:end

% item: NEW-U08-A08
% mdparagraph:P161:begin
Give the complete present and preterite indicative of ætēowian, then the singular and plural present and preterite subjunctive. State which persons each subjunctive entry covers.
% mdparagraph:P161:end

\end{minipage}\par\addvspace{4pt}

\par\begin{minipage}{\linewidth}\raggedright
% mdparagraph:P162:begin
\textbf{A9.}
% mdparagraph:P162:end

% item: NEW-U08-A09
% mdparagraph:P163:begin
Give the singular and plural present and preterite subjunctive of scīnan, stating which persons each entry covers. Then give both second-person commands of geblētsian and identify the number of people addressed.
% mdparagraph:P163:end

\end{minipage}\par\addvspace{4pt}

\reserve{7\baselineskip}
% mdparagraph:P164:begin
\subsection*{Practice Q. Questions}
% mdparagraph:P164:end

\reserve{5\baselineskip}
\par\begin{minipage}{\linewidth}\raggedright
% mdparagraph:P165:begin
\textbf{Q1.}
% mdparagraph:P165:end

% item: NEW-U08-Q01
% mdparagraph:P166:begin
Turn each statement into a yes/\allowbreak{}no question: Hē hæfð scip. Ðā męnn sindon gōde. Ðū gesāwe ðæt scip. Preserve person, number and tense.
% mdparagraph:P166:end

\end{minipage}\par\addvspace{4pt}

\par\begin{minipage}{\linewidth}\raggedright
% mdparagraph:P167:begin
\textbf{Q2.}
% mdparagraph:P167:end

% item: NEW-U08-Q02
% mdparagraph:P168:begin
Translate into Old English: Who has that ship? Whom does the man see? What does the man have? State the case of each interrogative.
% mdparagraph:P168:end

\end{minipage}\par\addvspace{4pt}

\par\begin{minipage}{\linewidth}\raggedright
% mdparagraph:P169:begin
\textbf{Q3.}
% mdparagraph:P169:end

% item: NEW-U08-Q03
% mdparagraph:P170:begin
Translate into Old English: To whom do you (sing.) give that fish? Whose ship do you (sing.) see? Explain what governs each interrogative case.
% mdparagraph:P170:end

\end{minipage}\par\addvspace{4pt}

\par\begin{minipage}{\linewidth}\raggedright
% mdparagraph:P171:begin
\textbf{Q4.}
% mdparagraph:P171:end

% item: NEW-U08-Q04
% mdparagraph:P172:begin
Two men are being compared. Ask “Which of the two men has that ship?” using attributive hwæðer with singular mann. Then ask “What sort of ship is that?” using hūlic.
% mdparagraph:P172:end

\end{minipage}\par\addvspace{4pt}

\par\begin{minipage}{\linewidth}\raggedright
% mdparagraph:P173:begin
\textbf{Q5.}
% mdparagraph:P173:end

% item: NEW-U08-Q05
% mdparagraph:P174:begin
Translate Hū swimð se mann? and Hwȳ swimð se mann? Explain the difference in the information requested.
% mdparagraph:P174:end

\end{minipage}\par\addvspace{4pt}

\par\begin{minipage}{\linewidth}\raggedright
% mdparagraph:P175:begin
\textbf{Q6.}
% mdparagraph:P175:end

% item: NEW-U08-Q06
% mdparagraph:P176:begin
Ask “Do you (sing.) have any fish?” using ǣnig and singular fisc. A learner adds dōn because English uses do; explain why that addition is unnecessary.
% mdparagraph:P176:end

\end{minipage}\par\addvspace{4pt}

\reserve{7\baselineskip}
% mdparagraph:P177:begin
\subsection*{Practice N. Pronouns and agreement}
% mdparagraph:P177:end

\reserve{5\baselineskip}
\par\begin{minipage}{\linewidth}\raggedright
% mdparagraph:P178:begin
\textbf{N1.}
% mdparagraph:P178:end

% item: NEW-U08-N01
% mdparagraph:P179:begin
Translate Sum mann hæfð scip and Sume męnn habbað scipu. Identify the subject and object forms and explain the change in number.
% mdparagraph:P179:end

\end{minipage}\par\addvspace{4pt}

\par\begin{minipage}{\linewidth}\raggedright
% mdparagraph:P180:begin
\textbf{N2.}
% mdparagraph:P180:end

% item: NEW-U08-N02
% mdparagraph:P181:begin
Translate: Each man sees a ship. Use gehwilc and singular mann. Then replace “each man” with “each companion” without changing the meaning of the verb.
% mdparagraph:P181:end

\end{minipage}\par\addvspace{4pt}

\par\begin{minipage}{\linewidth}\raggedright
% mdparagraph:P182:begin
\textbf{N3.}
% mdparagraph:P182:end

% item: NEW-U08-N03
% mdparagraph:P183:begin
Translate: The other man has another ship. Use ōðer twice. Explain why “other/\allowbreak{}another” is not automatically “second” here.
% mdparagraph:P183:end

\end{minipage}\par\addvspace{4pt}

\par\begin{minipage}{\linewidth}\raggedright
% mdparagraph:P184:begin
\textbf{N4.}
% mdparagraph:P184:end

% item: NEW-U08-N04
% mdparagraph:P185:begin
Translate: I saw the same man. I give the fish to the same man. Use ilca and give the case of each complete noun phrase.
% mdparagraph:P185:end

\end{minipage}\par\addvspace{4pt}

\par\begin{minipage}{\linewidth}\raggedright
% mdparagraph:P186:begin
\textbf{N5.}
% mdparagraph:P186:end

% item: NEW-U08-N05
% mdparagraph:P187:begin
Translate Hē self geseah ðæt scip. Then translate “I save myself,” spoken by a man, using mē selfne. Distinguish emphasis on the subject from a reflexive object.
% mdparagraph:P187:end

\end{minipage}\par\addvspace{4pt}

\par\begin{minipage}{\linewidth}\raggedright
% mdparagraph:P188:begin
\textbf{N6.}
% mdparagraph:P188:end

% item: NEW-U08-N06
% mdparagraph:P189:begin
Translate Swilce męnn sindon gōde. Explain how the inflected word swilce differs from the invariant adverb swilce “likewise.”
% mdparagraph:P189:end

\end{minipage}\par\addvspace{4pt}

\par\begin{minipage}{\linewidth}\raggedright
% mdparagraph:P190:begin
\textbf{N7.}
% mdparagraph:P190:end

% item: NEW-U08-N07
% mdparagraph:P191:begin
Give a full Old English noun phrase for each: of some men (genitive); to any companion (dative); that same ship (accusative); such little birds (accusative). Label the case and relevant agreement.
% mdparagraph:P191:end

\end{minipage}\par\addvspace{4pt}

\reserve{7\baselineskip}
% mdparagraph:P192:begin
\subsection*{Practice L. Relative clauses}
% mdparagraph:P192:end

\reserve{5\baselineskip}
\par\begin{minipage}{\linewidth}\raggedright
% mdparagraph:P193:begin
\textbf{L1.}
% mdparagraph:P193:end

% item: NEW-U08-L01
% mdparagraph:P194:begin
Join Se mann is gōd and Se mann hæfð ðæt scip into one sentence with ðe meaning “The man who has that ship is good.” Identify the relative subject.
% mdparagraph:P194:end

\end{minipage}\par\addvspace{4pt}

\par\begin{minipage}{\linewidth}\raggedright
% mdparagraph:P195:begin
\textbf{L2.}
% mdparagraph:P195:end

% item: NEW-U08-L02
% mdparagraph:P196:begin
Join Ic geseah ðæt scip and Se mann hæfð ðæt scip into one sentence with ðe. Retain “I saw” as the main clause.
% mdparagraph:P196:end

\end{minipage}\par\addvspace{4pt}

\par\begin{minipage}{\linewidth}\raggedright
% mdparagraph:P197:begin
\textbf{L3.}
% mdparagraph:P197:end

% item: NEW-U08-L03
% mdparagraph:P198:begin
Combine Ic sęlle ðone fisc ðǣm męnn and Se mann mē healp. Use an inflected demonstrative followed by ðe for the relative. Label the antecedent case and relative case separately.
% mdparagraph:P198:end

\end{minipage}\par\addvspace{4pt}

\par\begin{minipage}{\linewidth}\raggedright
% mdparagraph:P199:begin
\textbf{L4.}
% mdparagraph:P199:end

% item: NEW-U08-L04
% mdparagraph:P200:begin
Combine Ic geseah ðone mann and Ic healp ðǣm męnn. Use an inflected demonstrative followed by ðe. Explain why the two references to the man have different cases.
% mdparagraph:P200:end

\end{minipage}\par\addvspace{4pt}

\par\begin{minipage}{\linewidth}\raggedright
% mdparagraph:P201:begin
\textbf{L5.}
% mdparagraph:P201:end

% item: NEW-U08-L05
% mdparagraph:P202:begin
Translate: You (sing.), who helped me, give me that fish. Use ðe and the preterite of helpan; identify the relative verb’s person and the command’s person.
% mdparagraph:P202:end

\end{minipage}\par\addvspace{4pt}

\par\begin{minipage}{\linewidth}\raggedright
% mdparagraph:P203:begin
\textbf{L6.}
% mdparagraph:P203:end

% item: NEW-U08-L06
% mdparagraph:P204:begin
Change the addressee of your L5 sentence to two or more people. Change every affected pronoun and finite verb, keeping the fish singular.
% mdparagraph:P204:end

\end{minipage}\par\addvspace{4pt}

\par\begin{minipage}{\linewidth}\raggedright
% mdparagraph:P205:begin
\textbf{L7.}
% mdparagraph:P205:end

% item: NEW-U08-L07
% mdparagraph:P206:begin
In Se mann, ðe ðæt scip hæfð, is mīn gefēra and Ic geseah ðæt scip, ðe se mann hæfð, identify each antecedent and each relative role. Does ðe itself change its spelling to show those roles? Then translate Se mann, ðæs scip ic geseah, is gōd and explain the case and function of ðæs and scip.
% mdparagraph:P206:end

\end{minipage}\par\addvspace{4pt}

\par\begin{minipage}{\linewidth}\raggedright
% mdparagraph:P207:begin
\textbf{L8.}
% mdparagraph:P207:end

% item: NEW-U08-L08
% mdparagraph:P208:begin
A learner says that an inflected relative must always have the same case as its antecedent. Correct that rule using L3 or L4. State the ordinary production target and the limited exception noted in the lesson.
% mdparagraph:P208:end

\end{minipage}\par\addvspace{4pt}

\reserve{7\baselineskip}
% mdparagraph:P209:begin
\subsection*{Practice C. Lexical constructions}
% mdparagraph:P209:end

\reserve{5\baselineskip}
\par\begin{minipage}{\linewidth}\raggedright
% mdparagraph:P210:begin
\textbf{C1.}
% mdparagraph:P210:end

% item: NEW-U08-C01
% mdparagraph:P211:begin
Translate: They finish the work. They fill the waters. Use gefyllan in both sentences and explain the difference of sense.
% mdparagraph:P211:end

\end{minipage}\par\addvspace{4pt}

\par\begin{minipage}{\linewidth}\raggedright
% mdparagraph:P212:begin
\textbf{C2.}
% mdparagraph:P212:end

% item: NEW-U08-C02
% mdparagraph:P213:begin
Translate: The man has all the animals in his control. Use on with the appropriate form of geweald and an invariant possessor. Then make it a command to one person, “Have all the animals in your control.”
% mdparagraph:P213:end

\end{minipage}\par\addvspace{4pt}

\par\begin{minipage}{\linewidth}\raggedright
% mdparagraph:P214:begin
\textbf{C3.}
% mdparagraph:P214:end

% item: NEW-U08-C03
% mdparagraph:P215:begin
Translate: I live my life. We live our lives. State the case and number of each object and explain why libban can have that object here.
% mdparagraph:P215:end

\end{minipage}\par\addvspace{4pt}

\par\begin{minipage}{\linewidth}\raggedright
% mdparagraph:P216:begin
\textbf{C4.}
% mdparagraph:P216:end

% item: NEW-U08-C04
% mdparagraph:P217:begin
Translate: Dry land appears. Let dry land appear. Use ætēowian and identify the mood of each verb.
% mdparagraph:P217:end

\end{minipage}\par\addvspace{4pt}

\reserve{7\baselineskip}
% mdparagraph:P218:begin
\section*{8. Cook I.8}
% mdparagraph:P218:end

% mdparagraph:P219:begin
Translate into Old English. Cook’s complete exercise follows, with its original wording, parentheses and brackets. Where the English leaves an addressee’s number or the number of “fish” unspecified, a consistent defensible choice is possible. Retain an explicitly specified singular or plural.
% mdparagraph:P219:end

\reserve{5\baselineskip}
\par\begin{minipage}{\linewidth}\raggedright
% mdparagraph:P220:begin
\textbf{I.8.1.}
% mdparagraph:P220:end

% item: EX-I-08-01
% mdparagraph:P221:begin
Your works are good.
% mdparagraph:P221:end

\end{minipage}\par\addvspace{4pt}

\par\begin{minipage}{\linewidth}\raggedright
% mdparagraph:P222:begin
\textbf{I.8.2.}
% mdparagraph:P222:end

% item: EX-I-08-02
% mdparagraph:P223:begin
His tree is fruit-bearing.
% mdparagraph:P223:end

\end{minipage}\par\addvspace{4pt}

\par\begin{minipage}{\linewidth}\raggedright
% mdparagraph:P224:begin
\textbf{I.8.3.}
% mdparagraph:P224:end

% item: EX-I-08-03
% mdparagraph:P225:begin
Thou art the man.
% mdparagraph:P225:end

\end{minipage}\par\addvspace{4pt}

\par\begin{minipage}{\linewidth}\raggedright
% mdparagraph:P226:begin
\textbf{I.8.4.}
% mdparagraph:P226:end

% item: EX-I-08-04
% mdparagraph:P227:begin
This year is the fourth.
% mdparagraph:P227:end

\end{minipage}\par\addvspace{4pt}

\par\begin{minipage}{\linewidth}\raggedright
% mdparagraph:P228:begin
\textbf{I.8.5.}
% mdparagraph:P228:end

% item: EX-I-08-05
% mdparagraph:P229:begin
God made these little birds.
% mdparagraph:P229:end

\end{minipage}\par\addvspace{4pt}

\par\begin{minipage}{\linewidth}\raggedright
% mdparagraph:P230:begin
\textbf{I.8.6.}
% mdparagraph:P230:end

% item: EX-I-08-06
% mdparagraph:P231:begin
He created me likewise.
% mdparagraph:P231:end

\end{minipage}\par\addvspace{4pt}

\par\begin{minipage}{\linewidth}\raggedright
% mdparagraph:P232:begin
\textbf{I.8.7.}
% mdparagraph:P232:end

% item: EX-I-08-07
% mdparagraph:P233:begin
Our cattle (plur.) see the grass.
% mdparagraph:P233:end

\end{minipage}\par\addvspace{4pt}

\par\begin{minipage}{\linewidth}\raggedright
% mdparagraph:P234:begin
\textbf{I.8.8.}
% mdparagraph:P234:end

% item: EX-I-08-08
% mdparagraph:P235:begin
We live our lives (168.1).
% mdparagraph:P235:end

\end{minipage}\par\addvspace{4pt}

\par\begin{minipage}{\linewidth}\raggedright
% mdparagraph:P236:begin
\textbf{I.8.9.}
% mdparagraph:P236:end

% item: EX-I-08-09
% mdparagraph:P237:begin
His seeds grow in the earth.
% mdparagraph:P237:end

\end{minipage}\par\addvspace{4pt}

\par\begin{minipage}{\linewidth}\raggedright
% mdparagraph:P238:begin
\textbf{I.8.10.}
% mdparagraph:P238:end

% item: EX-I-08-10
% mdparagraph:P239:begin
The animals rest (themselves).
% mdparagraph:P239:end

\end{minipage}\par\addvspace{4pt}

\par\begin{minipage}{\linewidth}\raggedright
% mdparagraph:P240:begin
\textbf{I.8.11.}
% mdparagraph:P240:end

% item: EX-I-08-11
% mdparagraph:P241:begin
Fill the waters.
% mdparagraph:P241:end

\end{minipage}\par\addvspace{4pt}

\par\begin{minipage}{\linewidth}\raggedright
% mdparagraph:P242:begin
\textbf{I.8.12.}
% mdparagraph:P242:end

% item: EX-I-08-12
% mdparagraph:P243:begin
Let dry land appear.
% mdparagraph:P243:end

\end{minipage}\par\addvspace{4pt}

\par\begin{minipage}{\linewidth}\raggedright
% mdparagraph:P244:begin
\textbf{I.8.13.}
% mdparagraph:P244:end

% item: EX-I-08-13
% mdparagraph:P245:begin
Place (gesęte) the fish in the sea.
% mdparagraph:P245:end

\end{minipage}\par\addvspace{4pt}

% teaching: NEW-U08-COOK13-SUPPORT
% teaching: NEW-U08-VOC-gesettan
% teaching: NEW-U08-GRAM-u08-supplied-placement
% mdparagraph:P246:begin
\textbf{Supplied help for I.8.13.} Gesęte is the exact supplied singular command of weak-I gesęttan “set, place”. Use it as printed. For this placement, the thing placed is accusative and the location uses \textbf{on + dative}. Cook’s comparable placement in heaven has a dative location. This local help supplies the prefixed command and this bounded frame; its unaided recall is not required.
% mdparagraph:P246:end

\reserve{5\baselineskip}
\par\begin{minipage}{\linewidth}\raggedright
% mdparagraph:P247:begin
\textbf{I.8.14.}
% mdparagraph:P247:end

% item: EX-I-08-14
% mdparagraph:P248:begin
Have (sing.) all spirits under [in] your control.
% mdparagraph:P248:end

\end{minipage}\par\addvspace{4pt}

\par\begin{minipage}{\linewidth}\raggedright
% mdparagraph:P249:begin
\textbf{I.8.15.}
% mdparagraph:P249:end

% item: EX-I-08-15
% mdparagraph:P250:begin
Let the stars shine in the heavens.
% mdparagraph:P250:end

\end{minipage}\par\addvspace{4pt}

\par\begin{minipage}{\linewidth}\raggedright
% mdparagraph:P251:begin
\textbf{I.8.16.}
% mdparagraph:P251:end

% item: EX-I-08-16
% mdparagraph:P252:begin
I say indeed, Bless all men.
% mdparagraph:P252:end

\end{minipage}\par\addvspace{4pt}

\par\begin{minipage}{\linewidth}\raggedright
% mdparagraph:P253:begin
\textbf{I.8.17.}
% mdparagraph:P253:end

% item: EX-I-08-17
% mdparagraph:P254:begin
Thou, who madest life, bless me.
% mdparagraph:P254:end

\end{minipage}\par\addvspace{4pt}

\reserve{7\baselineskip}
% mdparagraph:P255:begin
\section*{9. The ship and the fish}
% mdparagraph:P255:end

% mdparagraph:P256:begin
\textbf{New teaching dialogue.} Three men are present. The first speaker asks his companion about the third man.
% mdparagraph:P256:end

% composition: NEW-U08-READING-01
% mdparagraph:P257:begin
\par\begin{minipage}{\linewidth}\raggedright
\readingtext{On ðǣm strande sindon ðrīe męnn and ān scip. Sum mann cwið: “Hwā is se mann, ðe ðæt scip hæfð?” His gefēra cwið: “Se mann is mīn gefēra. Hē hæfð ðā fiscas, ðe wē sēoð.” “Hwǣm sęlst ðū ðone fisc?” “Ic sęlle hine ðǣm męnn, se ðe mē healp.” “Hwilc scip hæfð hē?” “Hē hæfð ðæt lȳtle scip, ðæt ðe ðū gesāwe.”}
\end{minipage}\par
% mdparagraph:P257:end

\reserve{7\baselineskip}
% mdparagraph:P258:begin
\subsection*{Practice D. The connected dialogue}
% mdparagraph:P258:end

\reserve{5\baselineskip}
\par\begin{minipage}{\linewidth}\raggedright
% mdparagraph:P259:begin
\textbf{D1.}
% mdparagraph:P259:end

% item: NEW-U08-D01
% mdparagraph:P260:begin
Translate the complete account “The ship and the fish.” Keep the speakers and the three men’s relationships clear.
% mdparagraph:P260:end

\end{minipage}\par\addvspace{4pt}

\par\begin{minipage}{\linewidth}\raggedright
% mdparagraph:P261:begin
\textbf{D2.}
% mdparagraph:P261:end

% item: NEW-U08-D02
% mdparagraph:P262:begin
In the account, list the direct questions. For each interrogative, give its role and the information requested.
% mdparagraph:P262:end

\end{minipage}\par\addvspace{4pt}

\par\begin{minipage}{\linewidth}\raggedright
% mdparagraph:P263:begin
\textbf{D3.}
% mdparagraph:P263:end

% item: NEW-U08-D03
% mdparagraph:P264:begin
In the answer beginning Ic sęlle, explain the roles and cases of hine, ðǣm męnn, se and mē. What identifies the man receiving the fish?
% mdparagraph:P264:end

\end{minipage}\par\addvspace{4pt}

\par\begin{minipage}{\linewidth}\raggedright
% mdparagraph:P265:begin
\textbf{D4.}
% mdparagraph:P265:end

% item: NEW-U08-D04
% mdparagraph:P266:begin
Rewrite the final answer so that the man has two little ships which you (sing.) saw. Preserve its relative construction and change all necessary forms.
% mdparagraph:P266:end

\end{minipage}\par\addvspace{4pt}

\reserve{7\baselineskip}
% mdparagraph:P267:begin
\section*{10. The merchant and his merchandise}
% mdparagraph:P267:end

% mdparagraph:P268:begin
This exact four-turn opening is from Cook’s \emph{First Book}, “Trades and Occupations”, printed pp. 129–130 /\allowbreak{} one-based PDF pp. 145–146. Page 129 has three unnumbered dialogue lines, counted editorially as 1–3; page 130 supplies printed lines 1–7. The original English speaker labels are retained.
% mdparagraph:P268:end

\par\begin{minipage}{\linewidth}\raggedright
% source: FB-TRADES-MERCHANT-OPENING
% mdparagraph:P269:begin
\readingtext{\textbf{Teacher.} Hwæt sægst ðū, mangere?}
% mdparagraph:P269:end

% mdparagraph:P270:begin
\readingtext{\textbf{Merchant.} Ic sęcge ðæt behēfe ic eom ge cyninge and ealdormannum, and weligum, and eallum folce.}
% mdparagraph:P270:end

% mdparagraph:P271:begin
\readingtext{\textbf{Teacher.} And hū?}
% mdparagraph:P271:end

% mdparagraph:P272:begin
\readingtext{\textbf{Merchant.} Ic āstīge mīn scip mid hlæstum mīnum, and rōwe ofer sǣlice dǣlas, and cīepe mīn ðing, and bycge ðing dēorwierðu, ðā on ðissum lande ne bēoð ācęnnede, and ic hit tōgelǣde ēow hider mid miclum plihte ofer sǣ, and hwīlum forlidennesse ic ðolie, mid lyre ealra ðinga mīnra, unēaðe cwic ætberstende.}
% mdparagraph:P272:end

\end{minipage}\par
% mdparagraph:P273:begin
The excerpt ends with the merchant’s escape from shipwreck. The following question, “Hwilc ðing gelǣdst ðū ūs?”, and its goods catalogue are outside this selection.
% mdparagraph:P273:end

\reserve{7\baselineskip}
% mdparagraph:P274:begin
\subsection*{Local help}
% mdparagraph:P274:end

% teaching: NEW-U08-SOURCE-HELPS
% mdparagraph:P275:begin
All twelve groups below are supplied for the source-reading tasks. They give the unfamiliar words and constructions in their actual forms. The stated noun gender and verb or adjective class identify the headword; they do not require an unprinted complete paradigm for independent use. The question words and ordinary relative relation are learned material within these supported clauses.
% mdparagraph:P275:end

\reserve{5\baselineskip}
\par\begin{minipage}{\linewidth}\raggedright
% mdparagraph:P276:begin
\textbf{H1. mangere}
% mdparagraph:P276:end

% support: U08-SRC-HELP01
% teaching: NEW-U08-VOC-mangere
% teaching: NEW-U08-GRAM-u08-source-vocative
% mdparagraph:P277:begin
\textbf{Mangere}, masculine strong noun in -e: “merchant”.
% mdparagraph:P277:end

% mdparagraph:P278:begin
mangere, masculine, merchant; the nominative form is used in addressing him.
% mdparagraph:P278:end

\end{minipage}\par\addvspace{4pt}

\par\begin{minipage}{\linewidth}\raggedright
% mdparagraph:P279:begin
\textbf{H2. ðæt behēfe ic eom; And hū?}
% mdparagraph:P279:end

% support: U08-SRC-HELP02
% teaching: NEW-U08-VOC-behefe
% teaching: NEW-U08-VOC-thaet-conjunction
% teaching: NEW-U08-GRAM-u08-source-content-clause
% teaching: NEW-U08-GRAM-u08-source-useful-dative
% teaching: NEW-U08-GRAM-u08-source-elliptical-how
% mdparagraph:P280:begin
\textbf{Behēfe}, strong adjective with nominative -e: “useful”. \textbf{Ðæt} here is the supplied content conjunction.
% mdparagraph:P280:end

% mdparagraph:P281:begin
I say that I am useful. This ðæt introduces a content clause; it is not the neuter demonstrative or a relative. behēfe means useful and takes the dative people to whom the merchant is useful. The next question And hū? abbreviates And how [are you useful]?; its omitted predicate is recovered from this answer. That exact ellipsis is supplied here, without teaching unrestricted question ellipsis.
% mdparagraph:P281:end

\end{minipage}\par\addvspace{4pt}

\par\begin{minipage}{\linewidth}\raggedright
% mdparagraph:P282:begin
\textbf{H3. ge cyninge and ealdormannum, and weligum, and eallum folce}
% mdparagraph:P282:end

% support: U08-SRC-HELP03
% teaching: NEW-U08-VOC-ge-conjunction
% teaching: NEW-U08-VOC-cyning
% teaching: NEW-U08-VOC-ealdormann
% teaching: NEW-U08-VOC-welig
% teaching: NEW-U08-VOC-folc
% teaching: NEW-U08-GRAM-u08-source-ge-and
% teaching: NEW-U08-GRAM-u08-source-substantival-weligum
% mdparagraph:P283:begin
\textbf{Cyning}, masculine a-stem: “king”; \textbf{ealdormann}, masculine compound of mann: “leader, noble”; \textbf{welig}, strong adjective: “wealthy”; \textbf{folc}, neuter a-stem: “people”. \textbf{Ge … and} is the supplied correlative “both … and”.
% mdparagraph:P283:end

% mdparagraph:P284:begin
ge ... and coordinates the recipients: both to the king and to nobles, and to wealthy people, and to all the people. cyninge is dative singular; ealdormannum and weligum are dative plural; folce is dative singular. weligum is an adjective used substantivally.
% mdparagraph:P284:end

\end{minipage}\par\addvspace{4pt}

\par\begin{minipage}{\linewidth}\raggedright
% mdparagraph:P285:begin
\textbf{H4. āstīge mīn scip mid hlæstum mīnum}
% mdparagraph:P285:end

% support: U08-SRC-HELP04
% teaching: NEW-U08-VOC-astigan
% teaching: NEW-U08-VOC-hlaest
% teaching: NEW-U08-VOC-mid
% teaching: NEW-U08-GRAM-u08-source-board-cargo
% mdparagraph:P286:begin
\textbf{Āstīgan}, strong I, here “go aboard, board”; \textbf{hlæst}, neuter a-stem, plural here: “wares, cargo”. \textbf{Mid} means “with” in the locally supplied phrases.
% mdparagraph:P286:end

% mdparagraph:P287:begin
I board my ship with my cargo. āstīge is first singular present; scip is accusative, the thing boarded. mid governs dative hlæstum mīnum here. hlæst is neuter and is plural in this cargo expression.
% mdparagraph:P287:end

\end{minipage}\par\addvspace{4pt}

\par\begin{minipage}{\linewidth}\raggedright
% mdparagraph:P288:begin
\textbf{H5. rōwe ofer sǣlice dǣlas}
% mdparagraph:P288:end

% support: U08-SRC-HELP05
% teaching: NEW-U08-VOC-rowan
% teaching: NEW-U08-VOC-saelic
% teaching: NEW-U08-VOC-dael
% teaching: NEW-U08-GRAM-u08-source-nautical-motion
% mdparagraph:P289:begin
\textbf{Rōwan}, reduplicating strong verb: “row”; \textbf{sǣlic}, strong adjective: “of the sea”; \textbf{dǣl}, masculine a-stem: “part”.
% mdparagraph:P289:end

% mdparagraph:P290:begin
I row over sea regions /\allowbreak{} across stretches of sea. rōwe is first singular present; sǣlice dǣlas is accusative masculine plural after spatial ofer.
% mdparagraph:P290:end

\end{minipage}\par\addvspace{4pt}

\par\begin{minipage}{\linewidth}\raggedright
% mdparagraph:P291:begin
\textbf{H6. cīepe mīn ðing, and bycge ðing dēorwierðu}
% mdparagraph:P291:end

% support: U08-SRC-HELP06
% teaching: NEW-U08-VOC-ciepan
% teaching: NEW-U08-VOC-bycgan
% teaching: NEW-U08-VOC-deorwierth
% teaching: NEW-U08-GRAM-u08-source-buy-sell
% mdparagraph:P292:begin
\textbf{Cīepan}, weak I: “sell”; \textbf{bycgan}, irregular weak I: “buy”; \textbf{dēorwierðe}, strong adjective with nominative -e: “precious, valuable”.
% mdparagraph:P292:end

% mdparagraph:P293:begin
I sell my goods and buy valuable things. cīepe and bycge are first singular presents. Both ðing phrases are accusative plural; mīn and dēorwierðu agree with neuter plural ðing.
% mdparagraph:P293:end

\end{minipage}\par\addvspace{4pt}

\par\begin{minipage}{\linewidth}\raggedright
% mdparagraph:P294:begin
\textbf{H7. ðā on ðissum lande ne bēoð ācęnnede}
% mdparagraph:P294:end

% support: U08-SRC-HELP07
% teaching: NEW-U08-VOC-land
% teaching: NEW-U08-VOC-acennan
% teaching: NEW-U08-GRAM-u08-source-produced-passive
% mdparagraph:P295:begin
\textbf{Land}, neuter a-stem: “land, country”; \textbf{ācęnnan}, weak I: “produce, bring forth”. The entire negative passive \textbf{ne bēoð ācęnnede} is supplied here as “are not produced”.
% mdparagraph:P295:end

% mdparagraph:P296:begin
which are not produced in this land. The target relative ðā is nominative neuter plural, the subject inside its clause; its antecedent ðing is accusative in the main clause. lande is dative singular, neuter land. ācęnnede is a plural predicate past participle, are produced; its exact form is supplied, without assigning the whole verb or a general participial declension.
% mdparagraph:P296:end

\end{minipage}\par\addvspace{4pt}

\par\begin{minipage}{\linewidth}\raggedright
% mdparagraph:P297:begin
\textbf{H8. ic hit tōgelǣde ēow hider}
% mdparagraph:P297:end

% support: U08-SRC-HELP08
% teaching: NEW-U08-VOC-togelaedan
% teaching: NEW-U08-VOC-hider
% teaching: NEW-U08-GRAM-u08-source-bring-goods
% teaching: NEW-U08-GRAM-u08-source-collective-hit
% mdparagraph:P298:begin
\textbf{Tōgelǣdan}, weak I: “bring”; \textbf{hider}, direction adverb: “hither, here”.
% mdparagraph:P298:end

% mdparagraph:P299:begin
I bring it/\allowbreak{}the goods here to you. tōgelǣde is first singular present. ēow is dative plural recipient. hit is the actual neuter singular object used collectively for the merchandise; translate its reference, not as a new unrelated object. hider expresses movement hither.
% mdparagraph:P299:end

\end{minipage}\par\addvspace{4pt}

\par\begin{minipage}{\linewidth}\raggedright
% mdparagraph:P300:begin
\textbf{H9. mid miclum plihte ofer sǣ}
% mdparagraph:P300:end

% support: U08-SRC-HELP09
% teaching: NEW-U08-VOC-pliht
% teaching: NEW-U08-GRAM-u08-source-risk
% mdparagraph:P301:begin
\textbf{Pliht}, masculine a-stem: “peril, risk”. Mid introduces the supplied circumstance of risk here.
% mdparagraph:P301:end

% mdparagraph:P302:begin
at great risk across the sea. mid has a supplied circumstance use here and governs dative masculine singular miclum plihte. In ofer sǣ, sǣ has the same singular form in the accusative and dative. Across the sea motivates an accusative extent interpretation in this context, but the written form does not prove that case. No unique accusative answer is required from this surface alone; the task concerns the meaning across the sea.
% mdparagraph:P302:end

\end{minipage}\par\addvspace{4pt}

\par\begin{minipage}{\linewidth}\raggedright
% mdparagraph:P303:begin
\textbf{H10. hwīlum forlidennesse ic ðolie}
% mdparagraph:P303:end

% support: U08-SRC-HELP10
% teaching: NEW-U08-VOC-hwilum
% teaching: NEW-U08-VOC-forlidennes
% teaching: NEW-U08-VOC-tholian
% teaching: NEW-U08-GRAM-u08-source-shipwreck
% mdparagraph:P304:begin
\textbf{Hwīlum}, adverb: “sometimes”; \textbf{forlidennes}, feminine -nes noun: “shipwreck”; \textbf{ðolian}, weak II: “endure, experience”.
% mdparagraph:P304:end

% mdparagraph:P305:begin
sometimes I suffer shipwreck. hwīlum is a frequency adverb; forlidennesse is accusative feminine singular, the object of first singular present ðolie.
% mdparagraph:P305:end

\end{minipage}\par\addvspace{4pt}

\par\begin{minipage}{\linewidth}\raggedright
% mdparagraph:P306:begin
\textbf{H11. mid lyre ealra ðinga mīnra}
% mdparagraph:P306:end

% support: U08-SRC-HELP11
% teaching: NEW-U08-VOC-lyre
% teaching: NEW-U08-GRAM-u08-source-loss-genitive
% mdparagraph:P307:begin
\textbf{Lyre}, masculine strong noun in -e: “loss”. Mid introduces the supplied circumstance of loss.
% mdparagraph:P307:end

% mdparagraph:P308:begin
with the loss of all my things. lyre is dative masculine singular after mid; ealra ðinga mīnra is a genitive plural dependent of loss.
% mdparagraph:P308:end

\end{minipage}\par\addvspace{4pt}

\par\begin{minipage}{\linewidth}\raggedright
% mdparagraph:P309:begin
\textbf{H12. unēaðe cwic ætberstende}
% mdparagraph:P309:end

% support: U08-SRC-HELP12
% teaching: NEW-U08-VOC-uneathe
% teaching: NEW-U08-VOC-cwic
% teaching: NEW-U08-VOC-aetberstan
% teaching: NEW-U08-GRAM-u08-source-escape-participle
% mdparagraph:P310:begin
\textbf{Unēaðe}, adverb: “with difficulty, barely”; \textbf{cwic}, regular strong adjective: “alive”; \textbf{ætberstan}, strong III: “escape”. Cwic is related to the earlier irregular source form cucu, but retains its own actual spelling and analysis here.
% mdparagraph:P310:end

% mdparagraph:P311:begin
barely escaping alive. unēaðe means with difficulty/\allowbreak{}barely; cwic is nominative masculine singular referring to the merchant; ætberstende is a nominative singular present participle expressing his accompanying circumstance. The construction and exact forms are supplied.
% mdparagraph:P311:end

% mdparagraph:P312:begin
\textbf{Source presentation.} The Old English words, spelling, quantities and punctuation are retained. Printed line wrapping is joined, and the original note cues are moved to the page-qualified list below. The speaker labels remain the source’s English labels. In the notes, typographical emphasis is represented as plain text and line-end word divisions are joined; the continuation of p. 129 note 5 is reunited with its beginning.
% mdparagraph:P312:end

\end{minipage}\par\addvspace{4pt}

\reserve{7\baselineskip}
% mdparagraph:P313:begin
\subsection*{Cook’s original notes}
% mdparagraph:P313:end

% mdparagraph:P314:begin
These fifteen notes belong to Cook’s edition. Their numbers and grammar references are his, not the section numbers of this course. The local help above supplies the interpretation needed for the reading. All questions within these preserved notes are historical source commentary, not additional course tasks. The long etymological discussion in p. 129 note 5 is preserved as historical source commentary; it is not a new linguistic claim.
% mdparagraph:P314:end

% source-note: U08-ORIG-129-01
% mdparagraph:P315:begin
\textbf{P. 129, note 1 (sægst).} See 123.
% mdparagraph:P315:end

% source-note: U08-ORIG-129-02
% mdparagraph:P316:begin
\textbf{P. 129, note 2 (mangere).} Lat. mercator. Other Old English terms for merchant are cīepa and cīepmann. From a collateral form of the latter, cēapmann, without umlaut, is derived Mod. Eng. chapman. How is chap- related to cheap? See the New English Dictionary (New Eng. Dict.) under these words.
% mdparagraph:P316:end

% source-note: U08-ORIG-129-03
% mdparagraph:P317:begin
\textbf{P. 129, note 3 (behēfe).} Lat. utilis. Cf. the Mod. Eng. noun behoof.
% mdparagraph:P317:end

% source-note: U08-ORIG-129-04
% mdparagraph:P318:begin
\textbf{P. 129, note 4 (ge).} ge . . . and = Lat. et . . . et.
% mdparagraph:P318:end

% source-note: U08-ORIG-129-05
% mdparagraph:P319:begin
\textbf{P. 129, note 5 (cyninge).} Carlyle (Sartor Resartus, Bk. 3, Chap. 7) has the following: “König (King), anciently Könning, means Ken-ning (Cunning), or which is the same thing, Can-ning. Ever must the Sovereign of Mankind be fitly entitled King.” On the other hand Gummere (Germanic Origins, p. 270): “At the head of the family we found, of course, the father; and at the head of the state we naturally look for the king. The word ‘king’ means the child or son of the tribe, its representative or even creation; man of race, man of rank. Gradually the king ceases to be regarded as a creation of his race; his ancestry is pushed back to the gods, and his right is quite above all sanctions of popular choice or approval.” Which of these views is confirmed by etymology?
% mdparagraph:P319:end

% mdparagraph:P320:begin
(This original note continues on p. 130.)
% mdparagraph:P320:end

% source-note: U08-ORIG-129-06
% mdparagraph:P321:begin
\textbf{P. 129, note 6 (ealdormannum).} Lat. ducibus.
% mdparagraph:P321:end

% source-note: U08-ORIG-130-01
% mdparagraph:P322:begin
\textbf{P. 130, note 1 (hlæstum).} Lat. mercibus.
% mdparagraph:P322:end

% source-note: U08-ORIG-130-02
% mdparagraph:P323:begin
\textbf{P. 130, note 2 (rōwe).} Lat. navigo.
% mdparagraph:P323:end

% source-note: U08-ORIG-130-03
% mdparagraph:P324:begin
\textbf{P. 130, note 3 (sǣlice).} Lat. marinas.
% mdparagraph:P324:end

% source-note: U08-ORIG-130-04
% mdparagraph:P325:begin
\textbf{P. 130, note 4 (dǣlas).} Lat. partes.
% mdparagraph:P325:end

% source-note: U08-ORIG-130-05
% mdparagraph:P326:begin
\textbf{P. 130, note 5 (cīepe).} Lat. vendo.
% mdparagraph:P326:end

% source-note: U08-ORIG-130-06
% mdparagraph:P327:begin
\textbf{P. 130, note 6 (ðing/\allowbreak{}dēorwierðu).} Lat. res pretiosas.
% mdparagraph:P327:end

% source-note: U08-ORIG-130-07
% mdparagraph:P328:begin
\textbf{P. 130, note 7 (tōgelǣde).} Lat. adduco.
% mdparagraph:P328:end

% source-note: U08-ORIG-130-08
% mdparagraph:P329:begin
\textbf{P. 130, note 8 (plihte).} Lat. periculo. Mod. Eng. form of pliht?
% mdparagraph:P329:end

% source-note: U08-ORIG-130-09
% mdparagraph:P330:begin
\textbf{P. 130, note 9 (unēaðe/\allowbreak{}cwic/\allowbreak{}ætberstende).} Lat. vix vivus evadens. Note the love for alliteration, even in the Latin.
% mdparagraph:P330:end

\reserve{7\baselineskip}
% mdparagraph:P331:begin
\subsection*{Practice R. The original dialogue}
% mdparagraph:P331:end

\reserve{5\baselineskip}
\par\begin{minipage}{\linewidth}\raggedright
% mdparagraph:P332:begin
\textbf{R1.}
% mdparagraph:P332:end

% item: NEW-U08-R01
% mdparagraph:P333:begin
Translate the complete Merchant and Teacher excerpt using the local helps. Preserve who speaks each question and answer.
% mdparagraph:P333:end

\end{minipage}\par\addvspace{4pt}

\par\begin{minipage}{\linewidth}\raggedright
% mdparagraph:P334:begin
\textbf{R2.}
% mdparagraph:P334:end

% item: NEW-U08-R02
% mdparagraph:P335:begin
Identify the direct questions in the excerpt and explain what each asks. Use the local help for the merchant’s form of address.
% mdparagraph:P335:end

\end{minipage}\par\addvspace{4pt}

\par\begin{minipage}{\linewidth}\raggedright
% mdparagraph:P336:begin
\textbf{R3.}
% mdparagraph:P336:end

% item: NEW-U08-R03
% mdparagraph:P337:begin
In ðā on ðissum lande ne bēoð ācęnnede, identify the antecedent, the role of ðā and the number of the passive auxiliary. Use the supplied participle meaning.
% mdparagraph:P337:end

\end{minipage}\par\addvspace{4pt}

\par\begin{minipage}{\linewidth}\raggedright
% mdparagraph:P338:begin
\textbf{R4.}
% mdparagraph:P338:end

% item: NEW-U08-R04
% mdparagraph:P339:begin
Using the local helps, explain how the merchant obtains his goods and where he takes them. Distinguish these statements from the dangers he encounters.
% mdparagraph:P339:end

\end{minipage}\par\addvspace{4pt}

\par\begin{minipage}{\linewidth}\raggedright
% mdparagraph:P340:begin
\textbf{R5.}
% mdparagraph:P340:end

% item: NEW-U08-R05
% mdparagraph:P341:begin
Identify two expressions for dangers in the final answer, and explain the contribution of the closing words unēaðe cwic ætberstende. Use the supplied meanings and construction notes.
% mdparagraph:P341:end

\end{minipage}\par\addvspace{4pt}

\reserve{7\baselineskip}
% mdparagraph:P342:begin
\section*{11. Cumulative practice}
% mdparagraph:P342:end

\reserve{7\baselineskip}
% mdparagraph:P343:begin
\subsection*{Practice X. Questions, reference and clause relations}
% mdparagraph:P343:end

\reserve{5\baselineskip}
\par\begin{minipage}{\linewidth}\raggedright
% mdparagraph:P344:begin
\textbf{X1.}
% mdparagraph:P344:end

% item: NEW-U08-X01
% mdparagraph:P345:begin
Translate: Which men have the ships that you (sing.) saw? Use hwilc and ðe. Label the interrogative, antecedent and relative roles.
% mdparagraph:P345:end

\end{minipage}\par\addvspace{4pt}

\par\begin{minipage}{\linewidth}\raggedright
% mdparagraph:P346:begin
\textbf{X2.}
% mdparagraph:P346:end

% item: NEW-U08-X02
% mdparagraph:P347:begin
Translate: I give another fish to the same man who helped me. Use ōðer, ilca and an inflected demonstrative followed by ðe.
% mdparagraph:P347:end

\end{minipage}\par\addvspace{4pt}

\par\begin{minipage}{\linewidth}\raggedright
% mdparagraph:P348:begin
\textbf{X3.}
% mdparagraph:P348:end

% item: NEW-U08-X03
% mdparagraph:P349:begin
Translate: Do you (sing.) see any little birds? Use ǣnig with plural fugol and the appropriate form of lȳtel.
% mdparagraph:P349:end

\end{minipage}\par\addvspace{4pt}

\par\begin{minipage}{\linewidth}\raggedright
% mdparagraph:P350:begin
\textbf{X4.}
% mdparagraph:P350:end

% item: NEW-U08-X04
% mdparagraph:P351:begin
Translate: Each companion himself saw that ship. Use gehwilc and strong self. Then make the whole statement a direct yes/\allowbreak{}no question.
% mdparagraph:P351:end

\end{minipage}\par\addvspace{4pt}

\par\begin{minipage}{\linewidth}\raggedright
% mdparagraph:P352:begin
\textbf{X5.}
% mdparagraph:P352:end

% item: NEW-U08-X05
% mdparagraph:P353:begin
Translate: The man who has the ship says, “Let dry land appear.” Use ðe and retain a direct quotation.
% mdparagraph:P353:end

\end{minipage}\par\addvspace{4pt}

\par\begin{minipage}{\linewidth}\raggedright
% mdparagraph:P354:begin
\textbf{X6.}
% mdparagraph:P354:end

% item: NEW-U08-X06
% mdparagraph:P355:begin
Explain the different functions of ðæt in ðæt scip and ðæt ðe ðū gesāwe in the account. Does their identical spelling make both conjunctions meaning “that”?
% mdparagraph:P355:end

\end{minipage}\par\addvspace{4pt}

\clearpage
% mdparagraph:P356:begin
\section*{Source and editorial notes}
% mdparagraph:P356:end

\begingroup\fontsize{10}{12.8}\selectfont
\setlength{\parskip}{1.5pt}
% mdparagraph:P357:begin
The principal sources are Albert S. Cook, \emph{A First Book in Old English: Grammar, Reader, Notes, and Vocabulary}, second edition, revised and enlarged (Ginn, 1895), and \emph{Exercises in Old English} (Ginn, 1895). The complete I.8 prompts are printed on Exercises p. 4 /\allowbreak{} PDF p. 12. The original reader is the opening of “The Merchant and his Merchandise” on First Book pp. 129–130 /\allowbreak{} PDF pp. 145–146. No historical author is claimed for the sixteen teaching examples or the new ship dialogue.
% mdparagraph:P357:end

% mdparagraph:P358:begin
The pronoun systems and relative strategies follow Cook §§86–89, with the strong and weak endings in §§57–60; the supporting checks include Wright §§463–464 and 468–471 and Fulk §§35 and 77. The selected tables are complete paradigms built from checked lexical stems and inflection rules, not claims that every cell is individually attested in this excerpt. The bounded recognition note on case attraction qualifies the ordinary internal-case rule.
% mdparagraph:P358:end

% mdparagraph:P359:begin
The additional systems use Cook’s weak-II account in §§118–119, strong-I scīnan in §102, and the subjunctive endings in §§93–95. Cook §168.1 supplies the cognate accusative of libban. The source appearance form is ætēowie on First Book p. 124 /\allowbreak{} PDF p. 140, with note 20 directing the reader to ætēowian; the dictionary entry is on p. 275 /\allowbreak{} PDF p. 291. The contrasting transitive ætīewan is a different entry. The placement witness is p. 125 /\allowbreak{} PDF p. 141, printed line 15; the have-in-control witness is p. 127 /\allowbreak{} PDF p. 143, printed line 9. These short witnesses support the particular frames taught or supplied here; they do not assign the surrounding Creation passages.
% mdparagraph:P359:end

% mdparagraph:P360:begin
Dictionary compound hyphens are omitted in the displayed headwords ealdormann, āstīgan, dēorwierðe, ācęnnan, tōgelǣdan, forlidennes, unēaðe, ætberstan, ætēowian and gesęttan. Their running source spellings and quantities remain unchanged. Cook’s hūlic is retained beside the separately documented hūlīc spelling. All PDF pages cited here are one-based: the selected First Book scan uses printed page +16 and the Exercises scan +8; Wright uses +20 and Fulk +12.
% mdparagraph:P360:end

\endgroup
\end{document}
